%% Generated by Sphinx.
\def\sphinxdocclass{jupyterBook}
\documentclass[letterpaper,10pt,english]{jupyterBook}
\ifdefined\pdfpxdimen
   \let\sphinxpxdimen\pdfpxdimen\else\newdimen\sphinxpxdimen
\fi \sphinxpxdimen=.75bp\relax
\ifdefined\pdfimageresolution
    \pdfimageresolution= \numexpr \dimexpr1in\relax/\sphinxpxdimen\relax
\fi
%% let collapsible pdf bookmarks panel have high depth per default
\PassOptionsToPackage{bookmarksdepth=5}{hyperref}
%% turn off hyperref patch of \index as sphinx.xdy xindy module takes care of
%% suitable \hyperpage mark-up, working around hyperref-xindy incompatibility
\PassOptionsToPackage{hyperindex=false}{hyperref}
%% memoir class requires extra handling
\makeatletter\@ifclassloaded{memoir}
{\ifdefined\memhyperindexfalse\memhyperindexfalse\fi}{}\makeatother

\PassOptionsToPackage{warn}{textcomp}

\catcode`^^^^00a0\active\protected\def^^^^00a0{\leavevmode\nobreak\ }
\usepackage{cmap}
\usepackage{fontspec}
\defaultfontfeatures[\rmfamily,\sffamily,\ttfamily]{}
\usepackage{amsmath,amssymb,amstext}
\usepackage{polyglossia}
\setmainlanguage{english}



\setmainfont{FreeSerif}[
  Extension      = .otf,
  UprightFont    = *,
  ItalicFont     = *Italic,
  BoldFont       = *Bold,
  BoldItalicFont = *BoldItalic
]
\setsansfont{FreeSans}[
  Extension      = .otf,
  UprightFont    = *,
  ItalicFont     = *Oblique,
  BoldFont       = *Bold,
  BoldItalicFont = *BoldOblique,
]
\setmonofont{FreeMono}[
  Extension      = .otf,
  UprightFont    = *,
  ItalicFont     = *Oblique,
  BoldFont       = *Bold,
  BoldItalicFont = *BoldOblique,
]



\usepackage[Bjarne]{fncychap}
\usepackage[,numfigreset=1,mathnumfig]{sphinx}

\fvset{fontsize=\small}
\usepackage{geometry}


% Include hyperref last.
\usepackage{hyperref}
% Fix anchor placement for figures with captions.
\usepackage{hypcap}% it must be loaded after hyperref.
% Set up styles of URL: it should be placed after hyperref.
\urlstyle{same}

\addto\captionsenglish{\renewcommand{\contentsname}{Portada}}

\usepackage{sphinxmessages}



        % Start of preamble defined in sphinx-jupyterbook-latex %
         \usepackage[Latin,Greek]{ucharclasses}
        \usepackage{unicode-math}
        % fixing title of the toc
        \addto\captionsenglish{\renewcommand{\contentsname}{Contents}}
        \hypersetup{
            pdfencoding=auto,
            psdextra
        }
        % End of preamble defined in sphinx-jupyterbook-latex %
        

\title{Información Mutua, Fundamentos Matemáticos y Aplicaciones en Ciencia de Datos}
\date{Feb 10, 2026}
\release{}
\author{Juan Luis Aguayo}
\newcommand{\sphinxlogo}{\vbox{}}
\renewcommand{\releasename}{}
\makeindex
\begin{document}

\pagestyle{empty}
\sphinxmaketitle
\pagestyle{plain}
\sphinxtableofcontents
\pagestyle{normal}
\phantomsection\label{\detokenize{intro::doc}}


\sphinxAtStartPar
Este libro presenta una formulación moderna de la información mutua desde una perspectiva matemática y aplicada a ciencia de datos.

\sphinxAtStartPar
Se estudian, por ejemplo:
\begin{itemize}
\item {} 
\sphinxAtStartPar
Entropías

\item {} 
\sphinxAtStartPar
Información mutua

\item {} 
\sphinxAtStartPar
Divergencia KL

\item {} 
\sphinxAtStartPar
Modelos multivariados

\item {} 
\sphinxAtStartPar
Aplicaciones en machine learning

\end{itemize}

\sphinxstepscope


\part{Portada}

\sphinxstepscope


\chapter{Resumen}
\label{\detokenize{00_frontmatter/abstract:resumen}}\label{\detokenize{00_frontmatter/abstract::doc}}
\sphinxAtStartPar
La información mutua constituye una de las herramientas fundamentales dentro de la teoría de la información moderna y desempeña un rol central en múltiples áreas de la ingeniería, la estadística y la ciencia de datos. Su capacidad para cuantificar dependencias estadísticas generales —más allá de relaciones lineales— la convierte en un instrumento clave para el análisis de datos complejos, el aprendizaje automático y la inferencia probabilística.

\sphinxAtStartPar
Este libro presenta un tratamiento sistemático de la información mutua desde una perspectiva matemática, computacional y aplicada. Se desarrollan los fundamentos teóricos clásicos, incluyendo entropía, entropía condicional e información mutua, junto con sus propiedades estructurales, formulaciones multivariadas y conexiones con divergencias de información tales como la divergencia de Kullback–Leibler.

\sphinxAtStartPar
Posteriormente, se introducen modelos multivariados basados en matrices de información mutua, extensiones condicionales y análisis en espacios de alta dimensión, proporcionando herramientas formales para el estudio de dependencias complejas en datos reales. Estas construcciones se integran con aplicaciones modernas en ciencia de datos, incluyendo selección de variables, agrupamiento basado en información y métodos contemporáneos en aprendizaje profundo.

\sphinxAtStartPar
El texto combina el formalismo matemático con ejemplos interpretativos y desarrollos aplicados, manteniendo un equilibrio entre teoría y práctica. Asimismo, se incluyen demostraciones formales, tablas avanzadas y material complementario destinado a facilitar tanto la comprensión conceptual como la implementación práctica.

\sphinxAtStartPar
Este libro busca establecer un puente entre la teoría clásica de la información y las metodologías contemporáneas de análisis de datos, proporcionando una base sólida para el estudio avanzado de dependencias estadísticas en sistemas complejos.

\sphinxAtStartPar
{[}\hyperlink{cite.05_appendices/references:id6}{Shannon, 1948}{]},{[}\hyperlink{cite.05_appendices/references:id2}{Cover and Thomas, 2006}{]},{[}\hyperlink{cite.05_appendices/references:id3}{MacKay, 2003}{]}.

\sphinxstepscope


\chapter{Contribuciones del Libro}
\label{\detokenize{00_frontmatter/contributions:contribuciones-del-libro}}\label{\detokenize{00_frontmatter/contributions::doc}}
\sphinxAtStartPar
Este libro presenta un tratamiento integrado de la información mutua que combina formalismo matemático riguroso, modelado probabilístico moderno y aplicaciones contemporáneas en ciencia de datos e inteligencia artificial. Las contribuciones principales se estructuran en los siguientes ejes:


\section{1. Formalización Matemática Sistemática}
\label{\detokenize{00_frontmatter/contributions:formalizacion-matematica-sistematica}}
\sphinxAtStartPar
Se desarrolla una exposición rigurosa de los conceptos fundamentales de la teoría de la información:
\begin{itemize}
\item {} 
\sphinxAtStartPar
Entropía y entropía conjunta

\item {} 
\sphinxAtStartPar
Entropía condicional

\item {} 
\sphinxAtStartPar
Información mutua y sus variantes

\item {} 
\sphinxAtStartPar
Divergencia de Kullback–Leibler

\item {} 
\sphinxAtStartPar
Dependencias estadísticas generales

\end{itemize}

\sphinxAtStartPar
El enfoque enfatiza definiciones formales, demostraciones matemáticas y propiedades estructurales fundamentales.


\section{2. Integración entre Teoría de la Información y Probabilidad}
\label{\detokenize{00_frontmatter/contributions:integracion-entre-teoria-de-la-informacion-y-probabilidad}}
\sphinxAtStartPar
El texto establece conexiones explícitas entre:
\begin{itemize}
\item {} 
\sphinxAtStartPar
teoría de la medida probabilística,

\item {} 
\sphinxAtStartPar
variables aleatorias multivariadas,

\item {} 
\sphinxAtStartPar
modelos estadísticos,

\item {} 
\sphinxAtStartPar
dependencia e independencia condicional.

\end{itemize}

\sphinxAtStartPar
Se presenta una formulación coherente que permite interpretar la información mutua como una medida general de dependencia estadística.


\section{3. Extensiones Multivariadas y Modelos de Alta Dimensión}
\label{\detokenize{00_frontmatter/contributions:extensiones-multivariadas-y-modelos-de-alta-dimension}}
\sphinxAtStartPar
Se introducen estructuras avanzadas orientadas al análisis de sistemas complejos:
\begin{itemize}
\item {} 
\sphinxAtStartPar
matrices de información mutua,

\item {} 
\sphinxAtStartPar
información mutua condicional multivariable,

\item {} 
\sphinxAtStartPar
análisis de dependencias en espacios de alta dimensión,

\item {} 
\sphinxAtStartPar
representación estructural de relaciones estadísticas.

\end{itemize}

\sphinxAtStartPar
Estas herramientas permiten extender el análisis más allá del caso bivariado clásico.


\section{4. Conexión con Ciencia de Datos Moderna}
\label{\detokenize{00_frontmatter/contributions:conexion-con-ciencia-de-datos-moderna}}
\sphinxAtStartPar
El libro incorpora aplicaciones contemporáneas que vinculan teoría y práctica:
\begin{itemize}
\item {} 
\sphinxAtStartPar
selección de características basada en información,

\item {} 
\sphinxAtStartPar
agrupamiento mediante medidas informacionales,

\item {} 
\sphinxAtStartPar
análisis de redundancia y relevancia de variables,

\item {} 
\sphinxAtStartPar
interpretación de modelos probabilísticos complejos.

\end{itemize}

\sphinxAtStartPar
Se enfatiza el rol de la información mutua como herramienta transversal en análisis de datos.


\section{5. Integración con Aprendizaje Automático y Aprendizaje Profundo}
\label{\detokenize{00_frontmatter/contributions:integracion-con-aprendizaje-automatico-y-aprendizaje-profundo}}
\sphinxAtStartPar
Se analizan aplicaciones modernas en inteligencia artificial, incluyendo:
\begin{itemize}
\item {} 
\sphinxAtStartPar
estimación de información mutua en redes neuronales,

\item {} 
\sphinxAtStartPar
regularización basada en información,

\item {} 
\sphinxAtStartPar
representaciones latentes informacionales,

\item {} 
\sphinxAtStartPar
aprendizaje de representaciones y codificación eficiente.

\end{itemize}

\sphinxAtStartPar
Esto posiciona la teoría clásica dentro del contexto contemporáneo del aprendizaje automático.


\section{6. Enfoque Pedagógico y Estructural}
\label{\detokenize{00_frontmatter/contributions:enfoque-pedagogico-y-estructural}}
\sphinxAtStartPar
El libro ha sido diseñado para facilitar múltiples niveles de lectura:
\begin{itemize}
\item {} 
\sphinxAtStartPar
desarrollo progresivo desde fundamentos hasta aplicaciones avanzadas,

\item {} 
\sphinxAtStartPar
demostraciones formales complementadas con interpretaciones,

\item {} 
\sphinxAtStartPar
ejemplos teóricos y computacionales,

\item {} 
\sphinxAtStartPar
anexos técnicos con demostraciones extendidas y tablas avanzadas.

\end{itemize}

\sphinxAtStartPar
Este enfoque permite su uso tanto en cursos de posgrado como en investigación aplicada.


\section{7. Contribución Integradora}
\label{\detokenize{00_frontmatter/contributions:contribucion-integradora}}
\sphinxAtStartPar
La contribución principal del libro consiste en establecer un puente estructurado entre:
\begin{itemize}
\item {} 
\sphinxAtStartPar
teoría clásica de la información,

\item {} 
\sphinxAtStartPar
estadística matemática,

\item {} 
\sphinxAtStartPar
análisis multivariado,

\item {} 
\sphinxAtStartPar
ciencia de datos contemporánea,

\item {} 
\sphinxAtStartPar
aprendizaje profundo moderno.

\end{itemize}

\sphinxAtStartPar
El resultado es un marco unificado para el estudio avanzado de dependencias estadísticas en sistemas complejos y conjuntos de datos de alta dimensión.

\sphinxstepscope


\chapter{Notación y Convenciones Matemáticas}
\label{\detokenize{00_frontmatter/notation:notacion-y-convenciones-matematicas}}\label{\detokenize{00_frontmatter/notation::doc}}
\sphinxAtStartPar
Este capítulo establece las convenciones matemáticas, probabilísticas e informacionales utilizadas a lo largo del libro. La notación ha sido diseñada para mantener coherencia con la literatura científica moderna en teoría de la información, estadística matemática y ciencia de datos.


\bigskip\hrule\bigskip



\section{1. Conjuntos y Espacios}
\label{\detokenize{00_frontmatter/notation:conjuntos-y-espacios}}\begin{itemize}
\item {} 
\sphinxAtStartPar
\(\mathbb{R}\) : conjunto de los números reales.

\item {} 
\sphinxAtStartPar
\(\mathbb{R}^n\) : espacio euclidiano de dimensión \(n\).

\item {} 
\sphinxAtStartPar
\(\mathcal{X}, \mathcal{Y}, \mathcal{Z}\) : espacios de valores de variables aleatorias.

\item {} 
\sphinxAtStartPar
\(\Omega\) : espacio muestral.

\item {} 
\sphinxAtStartPar
\(\mathcal{F}\) : \(\sigma\)\sphinxhyphen{}álgebra asociada al espacio de probabilidad.

\item {} 
\sphinxAtStartPar
\((\Omega, \mathcal{F}, \mathbb{P})\) : espacio de probabilidad.

\end{itemize}


\bigskip\hrule\bigskip



\section{2. Variables Aleatorias y Distribuciones}
\label{\detokenize{00_frontmatter/notation:variables-aleatorias-y-distribuciones}}\begin{itemize}
\item {} 
\sphinxAtStartPar
\(X, Y, Z\) : variables aleatorias.

\item {} 
\sphinxAtStartPar
\(\mathbf{X}\) : vector aleatorio multivariado.

\item {} 
\sphinxAtStartPar
\(x, y, z\) : realizaciones específicas de variables aleatorias.

\item {} 
\sphinxAtStartPar
\(p(x)\) : función de masa de probabilidad.

\item {} 
\sphinxAtStartPar
\(f(x)\) : función de densidad de probabilidad.

\item {} 
\sphinxAtStartPar
\(p(x,y)\) : distribución conjunta.

\item {} 
\sphinxAtStartPar
\(p(x|y)\) : distribución condicional.

\end{itemize}

\sphinxAtStartPar
Distribuciones marginales:
\begin{equation*}
\begin{split}
p(x) = \sum_y p(x,y)
\end{split}
\end{equation*}
\sphinxAtStartPar
o en el caso continuo,
\begin{equation*}
\begin{split}
f(x) = \int f(x,y)\,dy.
\end{split}
\end{equation*}

\bigskip\hrule\bigskip



\section{3. Operadores Probabilísticos}
\label{\detokenize{00_frontmatter/notation:operadores-probabilisticos}}\begin{itemize}
\item {} 
\sphinxAtStartPar
\(\mathbb{P}(A)\) : probabilidad de un evento \(A\).

\item {} 
\sphinxAtStartPar
\(\mathbb{E}[X]\) : esperanza matemática.

\item {} 
\sphinxAtStartPar
\(\mathrm{Var}(X)\) : varianza.

\item {} 
\sphinxAtStartPar
\(\mathrm{Cov}(X,Y)\) : covarianza.

\end{itemize}

\sphinxAtStartPar
Esperanza condicional:
\begin{equation*}
\begin{split}
\mathbb{E}[X|Y].
\end{split}
\end{equation*}

\bigskip\hrule\bigskip



\section{4. Entropía e Información}
\label{\detokenize{00_frontmatter/notation:entropia-e-informacion}}
\sphinxAtStartPar
Entropía de Shannon:
\begin{equation*}
\begin{split}
H(X) = -\sum_x p(x)\log p(x).
\end{split}
\end{equation*}
\sphinxAtStartPar
Entropía conjunta:
\begin{equation*}
\begin{split}
H(X,Y).
\end{split}
\end{equation*}
\sphinxAtStartPar
Entropía condicional:
\begin{equation*}
\begin{split}
H(X|Y).
\end{split}
\end{equation*}
\sphinxAtStartPar
Información mutua:
\begin{equation*}
\begin{split}
I(X;Y).
\end{split}
\end{equation*}
\sphinxAtStartPar
Información mutua condicional:
\begin{equation*}
\begin{split}
I(X;Y|Z).
\end{split}
\end{equation*}
\sphinxAtStartPar
Divergencia de Kullback–Leibler:
\begin{equation*}
\begin{split}
D_{KL}(P\|Q).
\end{split}
\end{equation*}

\bigskip\hrule\bigskip



\section{5. Vectores y Matrices Informacionales}
\label{\detokenize{00_frontmatter/notation:vectores-y-matrices-informacionales}}\begin{itemize}
\item {} 
\sphinxAtStartPar
\(\mathbf{X} = (X_1,\dots,X_n)\) : vector aleatorio.

\item {} 
\sphinxAtStartPar
\(\Sigma\) : matriz de covarianza.

\item {} 
\sphinxAtStartPar
\(\mathbf{M}\) : matriz de información mutua.

\item {} 
\sphinxAtStartPar
\(M_{ij}\) : información mutua entre \(X_i\) y \(X_j\).

\end{itemize}


\bigskip\hrule\bigskip



\section{6. Convenciones Tipográficas}
\label{\detokenize{00_frontmatter/notation:convenciones-tipograficas}}\begin{itemize}
\item {} 
\sphinxAtStartPar
Variables aleatorias: letras mayúsculas (\(X, Y, Z\)).

\item {} 
\sphinxAtStartPar
Observaciones: letras minúsculas (\(x, y, z\)).

\item {} 
\sphinxAtStartPar
Vectores: negrita (\(\mathbf{X}\)).

\item {} 
\sphinxAtStartPar
Conjuntos: caligráfico (\(\mathcal{X}\)).

\item {} 
\sphinxAtStartPar
Operadores matemáticos: roman upright (ej. \(\mathrm{Var}\)).

\end{itemize}


\bigskip\hrule\bigskip



\section{7. Notación para Aprendizaje Automático}
\label{\detokenize{00_frontmatter/notation:notacion-para-aprendizaje-automatico}}\begin{itemize}
\item {} 
\sphinxAtStartPar
\(\mathbf{X}\) : matriz de datos.

\item {} 
\sphinxAtStartPar
\(\mathbf{z}\) : representación latente.

\item {} 
\sphinxAtStartPar
\(\theta\) : parámetros de modelos.

\item {} 
\sphinxAtStartPar
\(\mathcal{L}(\theta)\) : función de pérdida.

\item {} 
\sphinxAtStartPar
\(\hat{\theta}\) : estimador.

\end{itemize}


\bigskip\hrule\bigskip



\section{8. Convenciones Logarítmicas}
\label{\detokenize{00_frontmatter/notation:convenciones-logaritmicas}}
\sphinxAtStartPar
A menos que se indique lo contrario:
\begin{itemize}
\item {} 
\sphinxAtStartPar
\(\log\) denota logaritmo base 2.

\item {} 
\sphinxAtStartPar
Las unidades de información se expresan en bits.

\end{itemize}


\bigskip\hrule\bigskip



\section{9. Comentarios Finales}
\label{\detokenize{00_frontmatter/notation:comentarios-finales}}
\sphinxAtStartPar
Las definiciones matemáticas detalladas y demostraciones formales se presentan en capítulos posteriores. La notación aquí introducida se utilizará consistentemente a lo largo del texto para garantizar precisión conceptual y coherencia estructural.

\sphinxstepscope


\part{Fundamentos Teóricos}

\sphinxstepscope


\chapter{Entropía de Shannon}
\label{\detokenize{01_theory/entropy:entropia-de-shannon}}\label{\detokenize{01_theory/entropy::doc}}
\sphinxAtStartPar
La entropía es una de las nociones centrales de la teoría de la información. Introducida por Claude Shannon en 1948, cuantifica la incertidumbre promedio asociada a una variable aleatoria y constituye la base matemática para el análisis de sistemas de comunicación, aprendizaje automático y modelos probabilísticos.


\bigskip\hrule\bigskip



\section{1. Motivación Conceptual}
\label{\detokenize{01_theory/entropy:motivacion-conceptual}}
\sphinxAtStartPar
Considérese una variable aleatoria discreta \(X\) con distribución de probabilidad \(p(x)\). Si los resultados son altamente predecibles, la incertidumbre es baja; si son impredecibles, la incertidumbre es alta.

\sphinxAtStartPar
La entropía mide formalmente esta incertidumbre promedio.

\sphinxAtStartPar
Ejemplos intuitivos:
\begin{itemize}
\item {} 
\sphinxAtStartPar
Un dado cargado tiene menor entropía que un dado justo.

\item {} 
\sphinxAtStartPar
Una variable determinista tiene entropía cero.

\item {} 
\sphinxAtStartPar
Una distribución uniforme maximiza la entropía.

\end{itemize}


\bigskip\hrule\bigskip



\section{2. Definición Formal (Caso Discreto)}
\label{\detokenize{01_theory/entropy:definicion-formal-caso-discreto}}
\sphinxAtStartPar
Sea \(X\) una variable aleatoria discreta con función de masa \(p(x)\). La entropía de Shannon se define como:
\begin{equation*}
\begin{split}
H(X) = -\sum_{x \in \mathcal{X}} p(x)\log p(x).
\end{split}
\end{equation*}
\sphinxAtStartPar
donde:
\begin{itemize}
\item {} 
\sphinxAtStartPar
\(\log\) es el logaritmo base 2.

\item {} 
\sphinxAtStartPar
La unidad es el bit.

\end{itemize}


\bigskip\hrule\bigskip



\section{3. Interpretación Informacional}
\label{\detokenize{01_theory/entropy:interpretacion-informacional}}
\sphinxAtStartPar
La cantidad:
\begin{equation*}
\begin{split}
-\log p(x)
\end{split}
\end{equation*}
\sphinxAtStartPar
representa la información asociada al evento \(x\). Por tanto, la entropía es el valor esperado de la información:
\begin{equation*}
\begin{split}
H(X) = \mathbb{E}[-\log p(X)].
\end{split}
\end{equation*}
\sphinxAtStartPar
Interpretación:
\begin{itemize}
\item {} 
\sphinxAtStartPar
Alta entropía → mayor incertidumbre promedio.

\item {} 
\sphinxAtStartPar
Baja entropía → mayor previsibilidad.

\end{itemize}


\bigskip\hrule\bigskip



\section{4. Caso Continuo}
\label{\detokenize{01_theory/entropy:caso-continuo}}
\sphinxAtStartPar
Para variables aleatorias continuas con densidad \(f(x)\):
\begin{equation*}
\begin{split}
h(X) = -\int f(x)\log f(x)\,dx.
\end{split}
\end{equation*}
\sphinxAtStartPar
Esta cantidad se denomina entropía diferencial.

\sphinxAtStartPar
Advertencias:
\begin{itemize}
\item {} 
\sphinxAtStartPar
Puede ser negativa.

\item {} 
\sphinxAtStartPar
No es invariante ante transformaciones lineales.

\end{itemize}


\bigskip\hrule\bigskip



\section{5. Propiedades Fundamentales}
\label{\detokenize{01_theory/entropy:propiedades-fundamentales}}

\subsection{5.1 No negatividad}
\label{\detokenize{01_theory/entropy:no-negatividad}}\begin{equation*}
\begin{split}
H(X) \ge 0.
\end{split}
\end{equation*}

\subsection{5.2 Máximo para distribuciones uniformes}
\label{\detokenize{01_theory/entropy:maximo-para-distribuciones-uniformes}}
\sphinxAtStartPar
Si \(|\mathcal{X}| = n\):
\begin{equation*}
\begin{split}
H(X) \le \log n.
\end{split}
\end{equation*}

\subsection{5.3 Invarianza ante permutaciones}
\label{\detokenize{01_theory/entropy:invarianza-ante-permutaciones}}
\sphinxAtStartPar
La entropía depende únicamente de la distribución, no del etiquetado de los estados.


\subsection{5.4 Aditividad para variables independientes}
\label{\detokenize{01_theory/entropy:aditividad-para-variables-independientes}}
\sphinxAtStartPar
Si \(X\) y \(Y\) son independientes:
\begin{equation*}
\begin{split}
H(X,Y) = H(X) + H(Y).
\end{split}
\end{equation*}

\bigskip\hrule\bigskip



\section{6. Ejemplo Computacional}
\label{\detokenize{01_theory/entropy:ejemplo-computacional}}
\sphinxAtStartPar
Sea una variable binaria:
\begin{equation*}
\begin{split}
P(X=0)=0.5,\quad P(X=1)=0.5.
\end{split}
\end{equation*}
\sphinxAtStartPar
Entonces:
\begin{equation*}
\begin{split}
H(X)= -0.5\log 0.5 -0.5\log 0.5 = 1\ \text{bit}.
\end{split}
\end{equation*}
\sphinxAtStartPar
Caso sesgado:
\begin{equation*}
\begin{split}
P(X=1)=0.9,\quad P(X=0)=0.1.
\end{split}
\end{equation*}\begin{equation*}
\begin{split}
H(X) \approx 0.469\ \text{bits}.
\end{split}
\end{equation*}
\sphinxAtStartPar
Interpretación:
\begin{itemize}
\item {} 
\sphinxAtStartPar
Mayor sesgo → menor entropía.

\item {} 
\sphinxAtStartPar
Mayor equilibrio → mayor incertidumbre.

\end{itemize}


\bigskip\hrule\bigskip



\section{7. Interpretación Geométrica}
\label{\detokenize{01_theory/entropy:interpretacion-geometrica}}
\sphinxAtStartPar
La entropía puede interpretarse como una medida de dispersión probabilística sobre el simplex de probabilidades.
\begin{itemize}
\item {} 
\sphinxAtStartPar
Distribuciones concentradas → cercanas a vértices.

\item {} 
\sphinxAtStartPar
Distribuciones uniformes → centro del simplex.

\end{itemize}

\sphinxAtStartPar
Esta perspectiva es útil en:
\begin{itemize}
\item {} 
\sphinxAtStartPar
optimización

\item {} 
\sphinxAtStartPar
aprendizaje profundo

\item {} 
\sphinxAtStartPar
teoría de códigos

\end{itemize}


\bigskip\hrule\bigskip



\section{8. Rol en Ciencia de Datos}
\label{\detokenize{01_theory/entropy:rol-en-ciencia-de-datos}}
\sphinxAtStartPar
Aplicaciones modernas incluyen:
\begin{itemize}
\item {} 
\sphinxAtStartPar
selección de características

\item {} 
\sphinxAtStartPar
árboles de decisión

\item {} 
\sphinxAtStartPar
regularización probabilística

\item {} 
\sphinxAtStartPar
aprendizaje no supervisado

\item {} 
\sphinxAtStartPar
modelos generativos

\end{itemize}

\sphinxAtStartPar
En aprendizaje automático:
\begin{equation*}
\begin{split}
\text{Ganancia de información} = H(Y) - H(Y|X).
\end{split}
\end{equation*}

\bigskip\hrule\bigskip



\section{9. Conexión con Capítulos Posteriores}
\label{\detokenize{01_theory/entropy:conexion-con-capitulos-posteriores}}
\sphinxAtStartPar
La entropía sirve como base para:
\begin{itemize}
\item {} 
\sphinxAtStartPar
entropía condicional

\item {} 
\sphinxAtStartPar
información mutua

\item {} 
\sphinxAtStartPar
divergencia KL

\item {} 
\sphinxAtStartPar
matrices de información

\item {} 
\sphinxAtStartPar
modelos multivariados

\end{itemize}

\sphinxAtStartPar
Estos conceptos serán desarrollados en los capítulos siguientes.


\bigskip\hrule\bigskip



\section{10. Comentarios Finales}
\label{\detokenize{01_theory/entropy:comentarios-finales}}
\sphinxAtStartPar
La entropía establece el marco cuantitativo fundamental para medir incertidumbre e información. Su formulación probabilística permite conectar estadística, teoría de la comunicación y aprendizaje automático dentro de una estructura matemática unificada.

\sphinxAtStartPar
{[}\hyperlink{cite.05_appendices/references:id6}{Shannon, 1948}{]}, {[}\hyperlink{cite.05_appendices/references:id3}{MacKay, 2003}{]}, {[}\hyperlink{cite.05_appendices/references:id2}{Cover and Thomas, 2006}{]}

\sphinxstepscope


\chapter{Entropía Conjunta y Entropía Condicional}
\label{\detokenize{01_theory/conditional_entropy:entropia-conjunta-y-entropia-condicional}}\label{\detokenize{01_theory/conditional_entropy::doc}}

\section{Introducción}
\label{\detokenize{01_theory/conditional_entropy:introduccion}}
\sphinxAtStartPar
En teoría de la información, la \sphinxstylestrong{entropía conjunta} y la \sphinxstylestrong{entropía condicional}
permiten cuantificar la incertidumbre total de sistemas multivariables y la
incertidumbre residual de una variable cuando otra es conocida.

\sphinxAtStartPar
Estas cantidades constituyen la base formal para definir medidas más avanzadas
como la información mutua, la información mutua condicional y diversas medidas
de dependencia estadística en modelos multivariados.


\bigskip\hrule\bigskip



\section{Entropía Conjunta}
\label{\detokenize{01_theory/conditional_entropy:entropia-conjunta}}

\subsection{Definición}
\label{\detokenize{01_theory/conditional_entropy:definicion}}
\sphinxAtStartPar
Sean \(X\) y \(Y\) variables aleatorias discretas con distribución conjunta
\(P(X,Y)\). La entropía conjunta se define como
\begin{equation*}
\begin{split}
H(X,Y) = - \sum_{x} \sum_{y} P(x,y)\,\log P(x,y).
\end{split}
\end{equation*}

\bigskip\hrule\bigskip



\subsection{Interpretación}
\label{\detokenize{01_theory/conditional_entropy:interpretacion}}
\sphinxAtStartPar
La entropía conjunta mide:
\begin{itemize}
\item {} 
\sphinxAtStartPar
La incertidumbre total del sistema formado por \((X,Y)\).

\item {} 
\sphinxAtStartPar
La cantidad promedio de información necesaria para describir simultáneamente
ambas variables.

\end{itemize}


\bigskip\hrule\bigskip



\subsection{Propiedades}
\label{\detokenize{01_theory/conditional_entropy:propiedades}}\begin{enumerate}
\sphinxsetlistlabels{\arabic}{enumi}{enumii}{}{.}%
\item {} 
\sphinxAtStartPar
\sphinxstylestrong{Simetría}

\end{enumerate}
\begin{equation*}
\begin{split}
H(X,Y)=H(Y,X)
\end{split}
\end{equation*}\begin{enumerate}
\sphinxsetlistlabels{\arabic}{enumi}{enumii}{}{.}%
\setcounter{enumi}{1}
\item {} 
\sphinxAtStartPar
\sphinxstylestrong{No negatividad}

\end{enumerate}
\begin{equation*}
\begin{split}
H(X,Y)\ge 0
\end{split}
\end{equation*}\begin{enumerate}
\sphinxsetlistlabels{\arabic}{enumi}{enumii}{}{.}%
\setcounter{enumi}{2}
\item {} 
\sphinxAtStartPar
\sphinxstylestrong{Independencia}

\end{enumerate}

\sphinxAtStartPar
Si \(X\) e \(Y\) son independientes:
\begin{equation*}
\begin{split}
H(X,Y)=H(X)+H(Y)
\end{split}
\end{equation*}

\bigskip\hrule\bigskip



\subsection{Ejemplo Resuelto}
\label{\detokenize{01_theory/conditional_entropy:ejemplo-resuelto}}
\sphinxAtStartPar
Considere la distribución:


\begin{savenotes}\sphinxattablestart
\centering
\begin{tabulary}{\linewidth}[t]{|T|T|T|}
\hline
\sphinxstyletheadfamily 
\sphinxAtStartPar
X
&\sphinxstyletheadfamily 
\sphinxAtStartPar
Y
&\sphinxstyletheadfamily 
\sphinxAtStartPar
P(X,Y)
\\
\hline
\sphinxAtStartPar
0
&
\sphinxAtStartPar
0
&
\sphinxAtStartPar
0.25
\\
\hline
\sphinxAtStartPar
0
&
\sphinxAtStartPar
1
&
\sphinxAtStartPar
0.25
\\
\hline
\sphinxAtStartPar
1
&
\sphinxAtStartPar
0
&
\sphinxAtStartPar
0.25
\\
\hline
\sphinxAtStartPar
1
&
\sphinxAtStartPar
1
&
\sphinxAtStartPar
0.25
\\
\hline
\end{tabulary}
\par
\sphinxattableend\end{savenotes}

\sphinxAtStartPar
Entonces
\begin{equation*}
\begin{split}
H(X,Y) = -4(0.25\log 0.25)
\end{split}
\end{equation*}\begin{equation*}
\begin{split}
= -\log 0.25
\end{split}
\end{equation*}\begin{equation*}
\begin{split}
= 2 \text{ bits}.
\end{split}
\end{equation*}
\sphinxAtStartPar
Interpretación: el sistema posee incertidumbre máxima para dos variables binarias
equiprobables e independientes.


\bigskip\hrule\bigskip



\section{Entropía Condicional}
\label{\detokenize{01_theory/conditional_entropy:entropia-condicional}}

\subsection{Definición}
\label{\detokenize{01_theory/conditional_entropy:id1}}
\sphinxAtStartPar
La entropía condicional de \(X\) dado \(Y\) es
\begin{equation*}
\begin{split}
H(X|Y) = - \sum_{x,y} P(x,y)\log P(x|y).
\end{split}
\end{equation*}
\sphinxAtStartPar
Equivalentemente
\begin{equation*}
\begin{split}
H(X|Y) = \sum_y P(y)\,H(X|Y=y).
\end{split}
\end{equation*}

\bigskip\hrule\bigskip



\subsection{Interpretación}
\label{\detokenize{01_theory/conditional_entropy:id2}}
\sphinxAtStartPar
Mide:
\begin{itemize}
\item {} 
\sphinxAtStartPar
La incertidumbre restante de \(X\) después de conocer \(Y\).

\item {} 
\sphinxAtStartPar
El promedio de incertidumbre de \(X\) dentro de cada subconjunto definido por
valores de \(Y\).

\end{itemize}


\bigskip\hrule\bigskip



\subsection{Casos Extremos}
\label{\detokenize{01_theory/conditional_entropy:casos-extremos}}\begin{enumerate}
\sphinxsetlistlabels{\arabic}{enumi}{enumii}{}{.}%
\item {} 
\sphinxAtStartPar
\sphinxstylestrong{Dependencia Determinística}

\end{enumerate}

\sphinxAtStartPar
Si \(X=f(Y)\)
\begin{equation*}
\begin{split}
H(X|Y)=0
\end{split}
\end{equation*}\begin{enumerate}
\sphinxsetlistlabels{\arabic}{enumi}{enumii}{}{.}%
\setcounter{enumi}{1}
\item {} 
\sphinxAtStartPar
\sphinxstylestrong{Independencia}

\end{enumerate}

\sphinxAtStartPar
Si \(X\) e \(Y\) son independientes
\begin{equation*}
\begin{split}
H(X|Y)=H(X)
\end{split}
\end{equation*}

\bigskip\hrule\bigskip



\subsection{Ejemplo Resuelto}
\label{\detokenize{01_theory/conditional_entropy:id3}}
\sphinxAtStartPar
Suponga la distribución:


\begin{savenotes}\sphinxattablestart
\centering
\begin{tabulary}{\linewidth}[t]{|T|T|T|}
\hline
\sphinxstyletheadfamily 
\sphinxAtStartPar
Clima \(X\)
&\sphinxstyletheadfamily 
\sphinxAtStartPar
Congestión \(Y\)
&\sphinxstyletheadfamily 
\sphinxAtStartPar
P(X,Y)
\\
\hline
\sphinxAtStartPar
Soleado
&
\sphinxAtStartPar
No
&
\sphinxAtStartPar
0.50
\\
\hline
\sphinxAtStartPar
Soleado
&
\sphinxAtStartPar
Sí
&
\sphinxAtStartPar
0.10
\\
\hline
\sphinxAtStartPar
Lluvioso
&
\sphinxAtStartPar
No
&
\sphinxAtStartPar
0.15
\\
\hline
\sphinxAtStartPar
Lluvioso
&
\sphinxAtStartPar
Sí
&
\sphinxAtStartPar
0.25
\\
\hline
\end{tabulary}
\par
\sphinxattableend\end{savenotes}

\sphinxAtStartPar
Marginal de \(Y\):
\begin{equation*}
\begin{split}
P(Y=No)=0.65,\quad P(Y=Sí)=0.35.
\end{split}
\end{equation*}
\sphinxAtStartPar
Probabilidades condicionales:
\begin{equation*}
\begin{split}
P(Soleado|No)=0.769,\quad P(Lluvioso|No)=0.231
\end{split}
\end{equation*}\begin{equation*}
\begin{split}
P(Soleado|Sí)=0.286,\quad P(Lluvioso|Sí)=0.714
\end{split}
\end{equation*}
\sphinxAtStartPar
Entropías condicionadas:
\begin{equation*}
\begin{split}
H(X|No)=0.78 \text{ bits}
\end{split}
\end{equation*}\begin{equation*}
\begin{split}
H(X|Sí)=0.86 \text{ bits}
\end{split}
\end{equation*}
\sphinxAtStartPar
Entropía condicional total:
\begin{equation*}
\begin{split}
H(X|Y)=0.65(0.78)+0.35(0.86)=0.81 \text{ bits}.
\end{split}
\end{equation*}
\sphinxAtStartPar
Interpretación: conocer la congestión reduce parcialmente la incertidumbre del
clima, pero no completamente.


\bigskip\hrule\bigskip



\section{Relación Fundamental (Regla de la Cadena)}
\label{\detokenize{01_theory/conditional_entropy:relacion-fundamental-regla-de-la-cadena}}
\sphinxAtStartPar
La relación central entre estas cantidades es
\begin{equation*}
\begin{split}
H(X,Y)=H(Y)+H(X|Y)
\end{split}
\end{equation*}
\sphinxAtStartPar
y también
\begin{equation*}
\begin{split}
H(X,Y)=H(X)+H(Y|X).
\end{split}
\end{equation*}

\bigskip\hrule\bigskip



\subsection{Demostración Esquemática}
\label{\detokenize{01_theory/conditional_entropy:demostracion-esquematica}}
\sphinxAtStartPar
Partiendo de
\begin{equation*}
\begin{split}
P(x,y)=P(x|y)P(y)
\end{split}
\end{equation*}
\sphinxAtStartPar
entonces
\begin{equation*}
\begin{split}
H(X,Y) = -\sum_{x,y} P(x,y)\log[P(x|y)P(y)].
\end{split}
\end{equation*}
\sphinxAtStartPar
Separando logaritmos
\begin{equation*}
\begin{split}
H(X,Y)=H(Y)+H(X|Y).
\end{split}
\end{equation*}

\bigskip\hrule\bigskip



\section{Interpretación Geométrica}
\label{\detokenize{01_theory/conditional_entropy:interpretacion-geometrica}}
\sphinxAtStartPar
En términos conceptuales:
\begin{itemize}
\item {} 
\sphinxAtStartPar
\(H(X)\) representa la incertidumbre total de \(X\).

\item {} 
\sphinxAtStartPar
\(H(X|Y)\) es la porción no explicada por \(Y\).

\item {} 
\sphinxAtStartPar
La diferencia \(H(X)-H(X|Y)\) corresponde a información compartida
(información mutua).

\end{itemize}


\bigskip\hrule\bigskip



\section{Aplicaciones}
\label{\detokenize{01_theory/conditional_entropy:aplicaciones}}\begin{itemize}
\item {} 
\sphinxAtStartPar
Modelado de dependencia estadística

\item {} 
\sphinxAtStartPar
Selección de características

\item {} 
\sphinxAtStartPar
Sistemas de comunicación

\item {} 
\sphinxAtStartPar
Aprendizaje profundo

\item {} 
\sphinxAtStartPar
Modelos probabilísticos multivariados

\end{itemize}


\bigskip\hrule\bigskip



\section{Ejercicios Propuestos}
\label{\detokenize{01_theory/conditional_entropy:ejercicios-propuestos}}\begin{enumerate}
\sphinxsetlistlabels{\arabic}{enumi}{enumii}{}{.}%
\item {} 
\sphinxAtStartPar
Demuestre que

\end{enumerate}
\begin{equation*}
\begin{split}
H(X|Y)\le H(X).
\end{split}
\end{equation*}\begin{enumerate}
\sphinxsetlistlabels{\arabic}{enumi}{enumii}{}{.}%
\setcounter{enumi}{1}
\item {} 
\sphinxAtStartPar
Calcule \(H(X,Y)\) y \(H(X|Y)\) para dos variables binarias correlacionadas
con probabilidad de coincidencia \(p\).

\item {} 
\sphinxAtStartPar
Verifique la regla de la cadena para un ejemplo numérico propio.

\item {} 
\sphinxAtStartPar
Analice cómo cambia \(H(X|Y)\) cuando aumenta la correlación entre variables.

\end{enumerate}


\bigskip\hrule\bigskip



\section{Conclusión}
\label{\detokenize{01_theory/conditional_entropy:conclusion}}
\sphinxAtStartPar
La entropía conjunta y la entropía condicional constituyen herramientas
fundamentales para descomponer la incertidumbre en sistemas multivariados.
Estas cantidades permiten comprender la estructura informacional entre
variables y preparan el camino para la formulación formal de la información
mutua, estudiada en el siguiente capítulo.

\sphinxAtStartPar
{[}\hyperlink{cite.05_appendices/references:id3}{MacKay, 2003}{]}, {[}\hyperlink{cite.05_appendices/references:id2}{Cover and Thomas, 2006}{]}

\sphinxstepscope


\chapter{Información Mutua}
\label{\detokenize{01_theory/mutual_information:informacion-mutua}}\label{\detokenize{01_theory/mutual_information::doc}}

\section{Introducción}
\label{\detokenize{01_theory/mutual_information:introduccion}}
\sphinxAtStartPar
La información mutua es una medida fundamental de dependencia estadística
entre variables aleatorias. Cuantifica la reducción promedio de incertidumbre
de una variable cuando se conoce otra.

\sphinxAtStartPar
A diferencia de la correlación lineal, la información mutua captura
dependencias no lineales y relaciones complejas entre variables,
constituyendo una herramienta central en teoría de la información,
aprendizaje automático, estadística multivariada y ciencia de datos.


\bigskip\hrule\bigskip



\section{Definición Formal}
\label{\detokenize{01_theory/mutual_information:definicion-formal}}
\sphinxAtStartPar
Sean \(X\) e \(Y\) variables aleatorias discretas con distribución conjunta
\(P(X,Y)\) y marginales \(P(X)\) y \(P(Y)\). La información mutua se define como
\begin{equation*}
\begin{split}
I(X;Y)=\sum_{x,y} P(x,y)\log\frac{P(x,y)}{P(x)P(y)}.
\end{split}
\end{equation*}

\bigskip\hrule\bigskip



\section{Interpretaciones Equivalentes}
\label{\detokenize{01_theory/mutual_information:interpretaciones-equivalentes}}
\sphinxAtStartPar
La información mutua puede escribirse en términos de entropías como:


\subsection{Forma Entrópica Principal}
\label{\detokenize{01_theory/mutual_information:forma-entropica-principal}}\begin{equation*}
\begin{split}
I(X;Y)=H(X)+H(Y)-H(X,Y).
\end{split}
\end{equation*}

\subsection{Reducción de Incertidumbre}
\label{\detokenize{01_theory/mutual_information:reduccion-de-incertidumbre}}\begin{equation*}
\begin{split}
I(X;Y)=H(X)-H(X|Y).
\end{split}
\end{equation*}\begin{equation*}
\begin{split}
I(X;Y)=H(Y)-H(Y|X).
\end{split}
\end{equation*}
\sphinxAtStartPar
Estas expresiones muestran que la información mutua mide la disminución
de incertidumbre al conocer la otra variable.


\bigskip\hrule\bigskip



\section{Interpretación como Divergencia KL}
\label{\detokenize{01_theory/mutual_information:interpretacion-como-divergencia-kl}}
\sphinxAtStartPar
La información mutua también puede expresarse como una divergencia
de Kullback–Leibler:
\begin{equation*}
\begin{split}
I(X;Y)=D_{KL}\big(P(X,Y)\,\|\,P(X)P(Y)\big).
\end{split}
\end{equation*}
\sphinxAtStartPar
Interpretación:
\begin{itemize}
\item {} 
\sphinxAtStartPar
Mide qué tan distinta es la distribución conjunta respecto a la
distribución que existiría bajo independencia.

\item {} 
\sphinxAtStartPar
Si las variables son independientes, la divergencia es cero.

\end{itemize}


\bigskip\hrule\bigskip



\section{Propiedades Fundamentales}
\label{\detokenize{01_theory/mutual_information:propiedades-fundamentales}}

\subsection{1. No negatividad}
\label{\detokenize{01_theory/mutual_information:no-negatividad}}\begin{equation*}
\begin{split}
I(X;Y)\ge 0.
\end{split}
\end{equation*}

\subsection{2. Simetría}
\label{\detokenize{01_theory/mutual_information:simetria}}\begin{equation*}
\begin{split}
I(X;Y)=I(Y;X).
\end{split}
\end{equation*}

\subsection{3. Independencia}
\label{\detokenize{01_theory/mutual_information:independencia}}
\sphinxAtStartPar
Si \(X\) e \(Y\) son independientes:
\begin{equation*}
\begin{split}
I(X;Y)=0.
\end{split}
\end{equation*}

\subsection{4. Cota Superior}
\label{\detokenize{01_theory/mutual_information:cota-superior}}\begin{equation*}
\begin{split}
I(X;Y)\le \min(H(X),H(Y)).
\end{split}
\end{equation*}

\subsection{5. Invarianza ante transformaciones biyectivas}
\label{\detokenize{01_theory/mutual_information:invarianza-ante-transformaciones-biyectivas}}
\sphinxAtStartPar
La información mutua no cambia bajo transformaciones determinísticas
invertibles de las variables.


\bigskip\hrule\bigskip



\section{Ejemplo Resuelto: Clima y Congestión Vehicular}
\label{\detokenize{01_theory/mutual_information:ejemplo-resuelto-clima-y-congestion-vehicular}}
\sphinxAtStartPar
Considere la distribución:


\begin{savenotes}\sphinxattablestart
\centering
\begin{tabulary}{\linewidth}[t]{|T|T|T|}
\hline
\sphinxstyletheadfamily 
\sphinxAtStartPar
Clima \(X\)
&\sphinxstyletheadfamily 
\sphinxAtStartPar
Congestión \(Y\)
&\sphinxstyletheadfamily 
\sphinxAtStartPar
P(X,Y)
\\
\hline
\sphinxAtStartPar
Soleado
&
\sphinxAtStartPar
No
&
\sphinxAtStartPar
0.50
\\
\hline
\sphinxAtStartPar
Soleado
&
\sphinxAtStartPar
Sí
&
\sphinxAtStartPar
0.10
\\
\hline
\sphinxAtStartPar
Lluvioso
&
\sphinxAtStartPar
No
&
\sphinxAtStartPar
0.15
\\
\hline
\sphinxAtStartPar
Lluvioso
&
\sphinxAtStartPar
Sí
&
\sphinxAtStartPar
0.25
\\
\hline
\end{tabulary}
\par
\sphinxattableend\end{savenotes}

\sphinxAtStartPar
Marginales:
\begin{equation*}
\begin{split}
P(X=Soleado)=0.60,\quad P(X=Lluvioso)=0.40
\end{split}
\end{equation*}\begin{equation*}
\begin{split}
P(Y=No)=0.65,\quad P(Y=Sí)=0.35
\end{split}
\end{equation*}
\sphinxAtStartPar
Entropías previamente obtenidas:
\begin{equation*}
\begin{split}
H(X)=0.97 \text{ bits}
\end{split}
\end{equation*}\begin{equation*}
\begin{split}
H(X|Y)=0.81 \text{ bits}
\end{split}
\end{equation*}
\sphinxAtStartPar
Información mutua:
\begin{equation*}
\begin{split}
I(X;Y)=H(X)-H(X|Y)
\end{split}
\end{equation*}\begin{equation*}
\begin{split}
I(X;Y)=0.97-0.81=0.16 \text{ bits}.
\end{split}
\end{equation*}
\sphinxAtStartPar
Interpretación:
\begin{itemize}
\item {} 
\sphinxAtStartPar
Existe dependencia moderada entre clima y congestión.

\item {} 
\sphinxAtStartPar
Conocer la congestión reduce ligeramente la incertidumbre del clima.

\end{itemize}


\bigskip\hrule\bigskip



\section{Información Mutua Normalizada}
\label{\detokenize{01_theory/mutual_information:informacion-mutua-normalizada}}
\sphinxAtStartPar
Para comparar dependencias entre distintos sistemas se utilizan versiones
normalizadas, por ejemplo:
\begin{equation*}
\begin{split}
NMI=\frac{I(X;Y)}{\sqrt{H(X)H(Y)}}.
\end{split}
\end{equation*}
\sphinxAtStartPar
Aplicaciones:
\begin{itemize}
\item {} 
\sphinxAtStartPar
clustering

\item {} 
\sphinxAtStartPar
evaluación de modelos

\item {} 
\sphinxAtStartPar
análisis de dependencia relativa

\end{itemize}


\bigskip\hrule\bigskip



\section{Extensión a Variables Continuas}
\label{\detokenize{01_theory/mutual_information:extension-a-variables-continuas}}
\sphinxAtStartPar
Para variables continuas con densidad conjunta \(p(x,y)\):
\begin{equation*}
\begin{split}
I(X;Y)=\int\int p(x,y)\log\frac{p(x,y)}{p(x)p(y)}dx\,dy.
\end{split}
\end{equation*}
\sphinxAtStartPar
Observación:
\begin{itemize}
\item {} 
\sphinxAtStartPar
se usa entropía diferencial

\item {} 
\sphinxAtStartPar
pueden aparecer valores negativos en entropía diferencial, pero
la información mutua permanece no negativa

\end{itemize}


\bigskip\hrule\bigskip



\section{Interpretación Geométrica Informacional}
\label{\detokenize{01_theory/mutual_information:interpretacion-geometrica-informacional}}
\sphinxAtStartPar
Conceptualmente:
\begin{itemize}
\item {} 
\sphinxAtStartPar
\(H(X)\) es la incertidumbre total de \(X\)

\item {} 
\sphinxAtStartPar
\(H(X|Y)\) es la parte no explicada por \(Y\)

\item {} 
\sphinxAtStartPar
la intersección corresponde a \(I(X;Y)\)

\end{itemize}

\sphinxAtStartPar
Esta interpretación se usa frecuentemente en diagramas tipo Venn
informacionales.


\bigskip\hrule\bigskip



\section{Aplicaciones}
\label{\detokenize{01_theory/mutual_information:aplicaciones}}

\subsection{Aprendizaje Automático}
\label{\detokenize{01_theory/mutual_information:aprendizaje-automatico}}\begin{itemize}
\item {} 
\sphinxAtStartPar
selección de características

\item {} 
\sphinxAtStartPar
reducción de dimensionalidad

\item {} 
\sphinxAtStartPar
aprendizaje profundo

\item {} 
\sphinxAtStartPar
representación latente

\end{itemize}


\subsection{Ciencia de Datos}
\label{\detokenize{01_theory/mutual_information:ciencia-de-datos}}\begin{itemize}
\item {} 
\sphinxAtStartPar
análisis de dependencia no lineal

\item {} 
\sphinxAtStartPar
clustering basado en información

\item {} 
\sphinxAtStartPar
análisis exploratorio

\end{itemize}


\subsection{Comunicaciones}
\label{\detokenize{01_theory/mutual_information:comunicaciones}}\begin{itemize}
\item {} 
\sphinxAtStartPar
capacidad de canal

\item {} 
\sphinxAtStartPar
codificación óptima

\item {} 
\sphinxAtStartPar
análisis de ruido

\end{itemize}


\subsection{Estadística Multivariada}
\label{\detokenize{01_theory/mutual_information:estadistica-multivariada}}\begin{itemize}
\item {} 
\sphinxAtStartPar
matrices de información mutua

\item {} 
\sphinxAtStartPar
análisis de redes probabilísticas

\item {} 
\sphinxAtStartPar
modelos gráficos

\end{itemize}


\bigskip\hrule\bigskip



\section{Relación con Capítulos Posteriores}
\label{\detokenize{01_theory/mutual_information:relacion-con-capitulos-posteriores}}
\sphinxAtStartPar
Este concepto será extendido en:
\begin{itemize}
\item {} 
\sphinxAtStartPar
Matriz de Información Mutua

\item {} 
\sphinxAtStartPar
Información Mutua Condicional

\item {} 
\sphinxAtStartPar
Modelos Multivariados de Alta Dimensión

\item {} 
\sphinxAtStartPar
Divergencia de Kullback–Leibler

\end{itemize}


\bigskip\hrule\bigskip



\section{Ejercicios Propuestos}
\label{\detokenize{01_theory/mutual_information:ejercicios-propuestos}}\begin{enumerate}
\sphinxsetlistlabels{\arabic}{enumi}{enumii}{}{.}%
\item {} 
\sphinxAtStartPar
Demuestre que \(I(X;Y)\ge 0\) usando la desigualdad de Gibbs.

\item {} 
\sphinxAtStartPar
Calcule la información mutua para dos variables binarias con
probabilidad de coincidencia \(p\).

\item {} 
\sphinxAtStartPar
Compare información mutua y correlación lineal en un ejemplo no lineal.

\item {} 
\sphinxAtStartPar
Demuestre que

\end{enumerate}
\begin{equation*}
\begin{split}
I(X;Y)=0
\end{split}
\end{equation*}
\sphinxAtStartPar
si y solo si las variables son independientes.
\begin{enumerate}
\sphinxsetlistlabels{\arabic}{enumi}{enumii}{}{.}%
\setcounter{enumi}{4}
\item {} 
\sphinxAtStartPar
Construya un ejemplo donde la correlación sea cero pero la
información mutua sea positiva.

\end{enumerate}


\bigskip\hrule\bigskip



\section{Conclusión}
\label{\detokenize{01_theory/mutual_information:conclusion}}
\sphinxAtStartPar
La información mutua proporciona una medida universal de dependencia
estadística que trasciende relaciones lineales y permite caracterizar
estructuras informacionales complejas. Su formulación basada en entropía
y divergencia KL la convierte en una herramienta central para el análisis
teórico y aplicado en sistemas multivariados modernos.

\sphinxAtStartPar
{[}\hyperlink{cite.05_appendices/references:id6}{Shannon, 1948}{]}, {[}\hyperlink{cite.05_appendices/references:id2}{Cover and Thomas, 2006}{]}

\sphinxstepscope


\part{Fundamentos Matemáticos}

\sphinxstepscope


\chapter{Definiciones Formales}
\label{\detokenize{02_mathematical_foundations/formal_definitions:definiciones-formales}}\label{\detokenize{02_mathematical_foundations/formal_definitions::doc}}

\section{Introducción}
\label{\detokenize{02_mathematical_foundations/formal_definitions:introduccion}}
\sphinxAtStartPar
En los capítulos anteriores se presentaron las definiciones clásicas
de entropía e información mutua para variables discretas. Sin embargo,
para desarrollar resultados rigurosos y extensiones multivariadas es
necesario establecer un marco formal basado en teoría de la medida
y probabilidad moderna.

\sphinxAtStartPar
Este capítulo introduce:
\begin{itemize}
\item {} 
\sphinxAtStartPar
espacios de probabilidad

\item {} 
\sphinxAtStartPar
variables aleatorias medibles

\item {} 
\sphinxAtStartPar
distribuciones conjuntas y marginales

\item {} 
\sphinxAtStartPar
definiciones generales de entropía

\item {} 
\sphinxAtStartPar
información mutua en espacios generales

\end{itemize}


\bigskip\hrule\bigskip



\section{Espacio de Probabilidad}
\label{\detokenize{02_mathematical_foundations/formal_definitions:espacio-de-probabilidad}}
\sphinxAtStartPar
Un espacio de probabilidad se define como la terna
\begin{equation*}
\begin{split}
(\Omega,\mathcal{F},P)
\end{split}
\end{equation*}
\sphinxAtStartPar
donde:
\begin{itemize}
\item {} 
\sphinxAtStartPar
\(\Omega\) es el espacio muestral

\item {} 
\sphinxAtStartPar
\(\mathcal{F}\) es una \(\sigma\)\sphinxhyphen{}álgebra de eventos

\item {} 
\sphinxAtStartPar
\(P:\mathcal{F}\to[0,1]\) es una medida de probabilidad

\end{itemize}

\sphinxAtStartPar
cumpliendo:
\begin{enumerate}
\sphinxsetlistlabels{\arabic}{enumi}{enumii}{}{.}%
\item {} 
\sphinxAtStartPar
\(P(\Omega)=1\)

\item {} 
\sphinxAtStartPar
\(P(A)\ge0\)

\item {} 
\sphinxAtStartPar
aditividad numerable

\end{enumerate}
\begin{equation*}
\begin{split}
P\left(\bigcup_i A_i\right)=\sum_i P(A_i)
\end{split}
\end{equation*}
\sphinxAtStartPar
para eventos disjuntos.


\bigskip\hrule\bigskip



\section{Variables Aleatorias}
\label{\detokenize{02_mathematical_foundations/formal_definitions:variables-aleatorias}}
\sphinxAtStartPar
Una variable aleatoria es una función medible
\begin{equation*}
\begin{split}
X:\Omega\rightarrow \mathbb{R}
\end{split}
\end{equation*}
\sphinxAtStartPar
tal que
\begin{equation*}
\begin{split}
X^{-1}(B)\in \mathcal{F}
\end{split}
\end{equation*}
\sphinxAtStartPar
para todo conjunto boreliano \(B\).

\sphinxAtStartPar
Tipos relevantes:
\begin{itemize}
\item {} 
\sphinxAtStartPar
discretas

\item {} 
\sphinxAtStartPar
continuas

\item {} 
\sphinxAtStartPar
mixtas

\item {} 
\sphinxAtStartPar
vectoriales

\item {} 
\sphinxAtStartPar
procesos estocásticos

\end{itemize}


\bigskip\hrule\bigskip



\section{Distribuciones de Probabilidad}
\label{\detokenize{02_mathematical_foundations/formal_definitions:distribuciones-de-probabilidad}}
\sphinxAtStartPar
Sea \((X,Y)\) un vector aleatorio.


\subsection{Distribución Conjunta}
\label{\detokenize{02_mathematical_foundations/formal_definitions:distribucion-conjunta}}\begin{equation*}
\begin{split}
P_{X,Y}(A\times B)=P(X\in A, Y\in B).
\end{split}
\end{equation*}

\subsection{Distribuciones Marginales}
\label{\detokenize{02_mathematical_foundations/formal_definitions:distribuciones-marginales}}\begin{equation*}
\begin{split}
P_X(A)=P(X\in A)
\end{split}
\end{equation*}\begin{equation*}
\begin{split}
P_Y(B)=P(Y\in B)
\end{split}
\end{equation*}

\subsection{Independencia}
\label{\detokenize{02_mathematical_foundations/formal_definitions:independencia}}
\sphinxAtStartPar
Las variables son independientes si
\begin{equation*}
\begin{split}
P_{X,Y}=P_X\otimes P_Y.
\end{split}
\end{equation*}

\bigskip\hrule\bigskip



\section{Entropía General}
\label{\detokenize{02_mathematical_foundations/formal_definitions:entropia-general}}

\subsection{Caso Discreto}
\label{\detokenize{02_mathematical_foundations/formal_definitions:caso-discreto}}\begin{equation*}
\begin{split}
H(X)=-\sum_x P(x)\log P(x).
\end{split}
\end{equation*}

\subsection{Caso Continuo (Entropía Diferencial)}
\label{\detokenize{02_mathematical_foundations/formal_definitions:caso-continuo-entropia-diferencial}}
\sphinxAtStartPar
Si \(X\) tiene densidad \(f(x)\):
\begin{equation*}
\begin{split}
h(X)=-\int f(x)\log f(x)\,dx.
\end{split}
\end{equation*}
\sphinxAtStartPar
Observaciones:
\begin{itemize}
\item {} 
\sphinxAtStartPar
puede ser negativa

\item {} 
\sphinxAtStartPar
depende de la escala

\item {} 
\sphinxAtStartPar
no es invariante ante transformaciones generales

\end{itemize}


\bigskip\hrule\bigskip



\section{Entropía Condicional}
\label{\detokenize{02_mathematical_foundations/formal_definitions:entropia-condicional}}

\subsection{Forma General}
\label{\detokenize{02_mathematical_foundations/formal_definitions:forma-general}}\begin{equation*}
\begin{split}
H(X|Y)=\mathbb{E}_{Y}[H(X|Y=y)].
\end{split}
\end{equation*}
\sphinxAtStartPar
Caso discreto:
\begin{equation*}
\begin{split}
H(X|Y)=-\sum_{x,y}P(x,y)\log P(x|y).
\end{split}
\end{equation*}

\bigskip\hrule\bigskip



\section{Entropía Conjunta}
\label{\detokenize{02_mathematical_foundations/formal_definitions:entropia-conjunta}}\begin{equation*}
\begin{split}
H(X,Y)=-\sum_{x,y}P(x,y)\log P(x,y).
\end{split}
\end{equation*}
\sphinxAtStartPar
Regla de la cadena:
\begin{equation*}
\begin{split}
H(X,Y)=H(X)+H(Y|X).
\end{split}
\end{equation*}

\bigskip\hrule\bigskip



\section{Información Mutua General}
\label{\detokenize{02_mathematical_foundations/formal_definitions:informacion-mutua-general}}
\sphinxAtStartPar
La definición más general se formula mediante divergencia KL:
\begin{equation*}
\begin{split}
I(X;Y)=D_{KL}(P_{X,Y}\,\|\,P_X\otimes P_Y).
\end{split}
\end{equation*}
\sphinxAtStartPar
Caso discreto:
\begin{equation*}
\begin{split}
I(X;Y)=\sum_{x,y}P(x,y)
\log\frac{P(x,y)}{P(x)P(y)}.
\end{split}
\end{equation*}
\sphinxAtStartPar
Caso continuo:
\begin{equation*}
\begin{split}
I(X;Y)=\int\int p(x,y)\log
\frac{p(x,y)}{p(x)p(y)}dxdy.
\end{split}
\end{equation*}

\bigskip\hrule\bigskip



\section{Información Mutua Condicional}
\label{\detokenize{02_mathematical_foundations/formal_definitions:informacion-mutua-condicional}}
\sphinxAtStartPar
Definición general:
\begin{equation*}
\begin{split}
I(X;Y|Z)=
\mathbb{E}_Z
\left[
D_{KL}(P_{X,Y|Z}\,\|\,P_{X|Z}P_{Y|Z})
\right].
\end{split}
\end{equation*}
\sphinxAtStartPar
Forma entrópica:
\begin{equation*}
\begin{split}
I(X;Y|Z)=H(X|Z)-H(X|Y,Z).
\end{split}
\end{equation*}

\bigskip\hrule\bigskip



\section{Información Mutua Multivariada}
\label{\detokenize{02_mathematical_foundations/formal_definitions:informacion-mutua-multivariada}}
\sphinxAtStartPar
Para un vector aleatorio
\begin{equation*}
\begin{split}
\mathbf{X}=(X_1,\dots,X_n)
\end{split}
\end{equation*}
\sphinxAtStartPar
la información total se define como
\begin{equation*}
\begin{split}
TC(\mathbf{X})=
D_{KL}(P_{\mathbf{X}}\,
\|\,\prod_i P_{X_i}).
\end{split}
\end{equation*}
\sphinxAtStartPar
Interpretación:
\begin{itemize}
\item {} 
\sphinxAtStartPar
mide dependencia global

\item {} 
\sphinxAtStartPar
base para modelos multivariados

\item {} 
\sphinxAtStartPar
fundamento de matrices de información mutua

\end{itemize}


\bigskip\hrule\bigskip



\section{Variables Aleatorias Vectoriales}
\label{\detokenize{02_mathematical_foundations/formal_definitions:variables-aleatorias-vectoriales}}
\sphinxAtStartPar
Sea
\begin{equation*}
\begin{split}
\mathbf{X}\in \mathbb{R}^d.
\end{split}
\end{equation*}
\sphinxAtStartPar
La entropía diferencial multivariada:
\begin{equation*}
\begin{split}
h(\mathbf{X})=-\int f(\mathbf{x})
\log f(\mathbf{x})\,d\mathbf{x}.
\end{split}
\end{equation*}
\sphinxAtStartPar
Información mutua vectorial:
\begin{equation*}
\begin{split}
I(\mathbf{X};\mathbf{Y})
=D_{KL}(P_{\mathbf{X},\mathbf{Y}}
\|P_{\mathbf{X}}P_{\mathbf{Y}}).
\end{split}
\end{equation*}

\bigskip\hrule\bigskip



\section{Independencia Condicional}
\label{\detokenize{02_mathematical_foundations/formal_definitions:independencia-condicional}}
\sphinxAtStartPar
Las variables \(X\) e \(Y\) son independientes dado \(Z\) si
\begin{equation*}
\begin{split}
P_{X,Y|Z}=P_{X|Z}P_{Y|Z}.
\end{split}
\end{equation*}
\sphinxAtStartPar
Equivalente a
\begin{equation*}
\begin{split}
I(X;Y|Z)=0.
\end{split}
\end{equation*}

\bigskip\hrule\bigskip



\section{Notación Funcional}
\label{\detokenize{02_mathematical_foundations/formal_definitions:notacion-funcional}}
\sphinxAtStartPar
Se utilizarán las siguientes convenciones:
\begin{itemize}
\item {} 
\sphinxAtStartPar
letras mayúsculas: variables aleatorias

\item {} 
\sphinxAtStartPar
minúsculas: realizaciones

\item {} 
\sphinxAtStartPar
negrita: vectores aleatorios

\item {} 
\sphinxAtStartPar
\(\mathbb{E}\): esperanza

\item {} 
\sphinxAtStartPar
\(D_{KL}\): divergencia de Kullback–Leibler

\end{itemize}


\bigskip\hrule\bigskip



\section{Ejemplo Formal Discreto}
\label{\detokenize{02_mathematical_foundations/formal_definitions:ejemplo-formal-discreto}}
\sphinxAtStartPar
Sea:
\begin{equation*}
\begin{split}
X,Y\in\{0,1\}.
\end{split}
\end{equation*}
\sphinxAtStartPar
Distribución conjunta:
\begin{equation*}
\begin{split}
P(0,0)=0.4,\quad
P(0,1)=0.1,
\end{split}
\end{equation*}\begin{equation*}
\begin{split}
P(1,0)=0.2,\quad
P(1,1)=0.3.
\end{split}
\end{equation*}
\sphinxAtStartPar
Marginales:
\begin{equation*}
\begin{split}
P_X(0)=0.5,\quad
P_X(1)=0.5
\end{split}
\end{equation*}\begin{equation*}
\begin{split}
P_Y(0)=0.6,\quad
P_Y(1)=0.4.
\end{split}
\end{equation*}
\sphinxAtStartPar
Información mutua:
\begin{equation*}
\begin{split}
I(X;Y)=\sum_{x,y}P(x,y)
\log\frac{P(x,y)}{P(x)P(y)}.
\end{split}
\end{equation*}
\sphinxAtStartPar
Este ejemplo se retomará en capítulos posteriores para:
\begin{itemize}
\item {} 
\sphinxAtStartPar
matrices de información mutua

\item {} 
\sphinxAtStartPar
modelos multivariados

\item {} 
\sphinxAtStartPar
clustering informacional

\end{itemize}


\bigskip\hrule\bigskip



\section{Conclusión}
\label{\detokenize{02_mathematical_foundations/formal_definitions:conclusion}}
\sphinxAtStartPar
Las definiciones formales introducidas permiten extender los conceptos
clásicos de entropía e información mutua hacia:
\begin{itemize}
\item {} 
\sphinxAtStartPar
espacios generales de probabilidad

\item {} 
\sphinxAtStartPar
variables continuas y vectoriales

\item {} 
\sphinxAtStartPar
modelos multivariados complejos

\item {} 
\sphinxAtStartPar
ciencia de datos moderna

\end{itemize}

\sphinxAtStartPar
Este formalismo servirá como base para el desarrollo riguroso de
propiedades matemáticas y teoremas en el siguiente capítulo.

\sphinxAtStartPar
{[}\hyperlink{cite.05_appendices/references:id10}{Csiszár, 2011}{]}, {[}\hyperlink{cite.05_appendices/references:id2}{Cover and Thomas, 2006}{]}

\sphinxstepscope


\chapter{Propiedades y Teoremas Fundamentales}
\label{\detokenize{02_mathematical_foundations/properties_theorems:propiedades-y-teoremas-fundamentales}}\label{\detokenize{02_mathematical_foundations/properties_theorems::doc}}

\section{Introducción}
\label{\detokenize{02_mathematical_foundations/properties_theorems:introduccion}}
\sphinxAtStartPar
En este capítulo se presentan las propiedades matemáticas esenciales
de la entropía, la entropía condicional y la información mutua.
Estos resultados constituyen la base teórica para:
\begin{itemize}
\item {} 
\sphinxAtStartPar
modelos multivariados

\item {} 
\sphinxAtStartPar
aprendizaje automático basado en información

\item {} 
\sphinxAtStartPar
selección de características

\item {} 
\sphinxAtStartPar
análisis de dependencia estadística

\end{itemize}

\sphinxAtStartPar
Se incluyen demostraciones formales y resultados clásicos de la
teoría de la información.


\bigskip\hrule\bigskip



\section{No Negatividad de la Entropía}
\label{\detokenize{02_mathematical_foundations/properties_theorems:no-negatividad-de-la-entropia}}

\subsection{Teorema}
\label{\detokenize{02_mathematical_foundations/properties_theorems:teorema}}
\sphinxAtStartPar
Para toda variable aleatoria discreta:
\begin{equation*}
\begin{split}
H(X)\ge 0.
\end{split}
\end{equation*}

\subsection{Demostración}
\label{\detokenize{02_mathematical_foundations/properties_theorems:demostracion}}
\sphinxAtStartPar
Dado que \(0\le P(x)\le 1\) se tiene:
\begin{equation*}
\begin{split}
\log P(x)\le 0.
\end{split}
\end{equation*}
\sphinxAtStartPar
Entonces:
\begin{equation*}
\begin{split}
- P(x)\log P(x)\ge 0.
\end{split}
\end{equation*}
\sphinxAtStartPar
Sumando sobre todos los valores:
\begin{equation*}
\begin{split}
H(X)\ge 0.
\end{split}
\end{equation*}

\bigskip\hrule\bigskip



\section{Entropía Máxima}
\label{\detokenize{02_mathematical_foundations/properties_theorems:entropia-maxima}}

\subsection{Teorema}
\label{\detokenize{02_mathematical_foundations/properties_theorems:id1}}
\sphinxAtStartPar
Entre todas las distribuciones discretas con soporte finito de tamaño \(n\),
la entropía es máxima para la distribución uniforme:
\begin{equation*}
\begin{split}
P(x)=\frac{1}{n}.
\end{split}
\end{equation*}

\subsection{Resultado}
\label{\detokenize{02_mathematical_foundations/properties_theorems:resultado}}\begin{equation*}
\begin{split}
H(X)\le \log n.
\end{split}
\end{equation*}
\sphinxAtStartPar
La igualdad ocurre solo en el caso equiprobable.


\bigskip\hrule\bigskip



\section{Regla de la Cadena}
\label{\detokenize{02_mathematical_foundations/properties_theorems:regla-de-la-cadena}}

\subsection{Teorema}
\label{\detokenize{02_mathematical_foundations/properties_theorems:id2}}\begin{equation*}
\begin{split}
H(X,Y)=H(X)+H(Y|X).
\end{split}
\end{equation*}

\subsection{Demostración}
\label{\detokenize{02_mathematical_foundations/properties_theorems:id3}}
\sphinxAtStartPar
Partimos de:
\begin{equation*}
\begin{split}
P(x,y)=P(x)P(y|x).
\end{split}
\end{equation*}
\sphinxAtStartPar
Luego:
\begin{equation*}
\begin{split}
\log P(x,y)=\log P(x)+\log P(y|x).
\end{split}
\end{equation*}
\sphinxAtStartPar
Sustituyendo en la definición de entropía conjunta:
\begin{equation*}
\begin{split}
H(X,Y)=-\sum P(x,y)\log P(x,y).
\end{split}
\end{equation*}
\sphinxAtStartPar
Se obtiene:
\begin{equation*}
\begin{split}
H(X,Y)=H(X)+H(Y|X).
\end{split}
\end{equation*}

\bigskip\hrule\bigskip



\section{Subaditividad}
\label{\detokenize{02_mathematical_foundations/properties_theorems:subaditividad}}

\subsection{Teorema}
\label{\detokenize{02_mathematical_foundations/properties_theorems:id4}}\begin{equation*}
\begin{split}
H(X,Y)\le H(X)+H(Y).
\end{split}
\end{equation*}

\subsection{Demostración}
\label{\detokenize{02_mathematical_foundations/properties_theorems:id5}}
\sphinxAtStartPar
Usando:
\begin{equation*}
\begin{split}
H(Y|X)\le H(Y).
\end{split}
\end{equation*}
\sphinxAtStartPar
y la regla de la cadena:
\begin{equation*}
\begin{split}
H(X,Y)=H(X)+H(Y|X).
\end{split}
\end{equation*}
\sphinxAtStartPar
se concluye la desigualdad.


\bigskip\hrule\bigskip



\section{Desigualdad de Gibbs}
\label{\detokenize{02_mathematical_foundations/properties_theorems:desigualdad-de-gibbs}}

\subsection{Teorema}
\label{\detokenize{02_mathematical_foundations/properties_theorems:id6}}
\sphinxAtStartPar
Para distribuciones \(P\) y \(Q\):
\begin{equation*}
\begin{split}
D_{KL}(P\|Q)\ge 0.
\end{split}
\end{equation*}

\subsection{Demostración (Jensen)}
\label{\detokenize{02_mathematical_foundations/properties_theorems:demostracion-jensen}}
\sphinxAtStartPar
La función \(-\log x\) es convexa.
Aplicando desigualdad de Jensen:
\begin{equation*}
\begin{split}
\mathbb{E}\left[\log\frac{Q}{P}\right]
\le
\log\mathbb{E}\left[\frac{Q}{P}\right].
\end{split}
\end{equation*}
\sphinxAtStartPar
Se concluye:
\begin{equation*}
\begin{split}
D_{KL}(P\|Q)\ge 0.
\end{split}
\end{equation*}

\bigskip\hrule\bigskip



\section{No Negatividad de la Información Mutua}
\label{\detokenize{02_mathematical_foundations/properties_theorems:no-negatividad-de-la-informacion-mutua}}

\subsection{Teorema}
\label{\detokenize{02_mathematical_foundations/properties_theorems:id7}}\begin{equation*}
\begin{split}
I(X;Y)\ge 0.
\end{split}
\end{equation*}

\subsection{Demostración}
\label{\detokenize{02_mathematical_foundations/properties_theorems:id8}}
\sphinxAtStartPar
Usando:
\begin{equation*}
\begin{split}
I(X;Y)=D_{KL}(P_{X,Y}\|P_XP_Y)
\end{split}
\end{equation*}
\sphinxAtStartPar
y la desigualdad de Gibbs.


\bigskip\hrule\bigskip



\section{Simetría}
\label{\detokenize{02_mathematical_foundations/properties_theorems:simetria}}

\subsection{Propiedad}
\label{\detokenize{02_mathematical_foundations/properties_theorems:propiedad}}\begin{equation*}
\begin{split}
I(X;Y)=I(Y;X).
\end{split}
\end{equation*}
\sphinxAtStartPar
Esto se deduce directamente de:
\begin{equation*}
\begin{split}
H(X)+H(Y)-H(X,Y).
\end{split}
\end{equation*}

\bigskip\hrule\bigskip



\section{Información Mutua Nula}
\label{\detokenize{02_mathematical_foundations/properties_theorems:informacion-mutua-nula}}

\subsection{Teorema}
\label{\detokenize{02_mathematical_foundations/properties_theorems:id9}}\begin{equation*}
\begin{split}
I(X;Y)=0 \iff X\perp Y.
\end{split}
\end{equation*}

\subsection{Interpretación}
\label{\detokenize{02_mathematical_foundations/properties_theorems:interpretacion}}
\sphinxAtStartPar
La información mutua mide dependencia estadística.


\bigskip\hrule\bigskip



\section{Regla de la Cadena para Información Mutua}
\label{\detokenize{02_mathematical_foundations/properties_theorems:regla-de-la-cadena-para-informacion-mutua}}

\subsection{Teorema}
\label{\detokenize{02_mathematical_foundations/properties_theorems:id10}}\begin{equation*}
\begin{split}
I(X;Y,Z)=I(X;Y)+I(X;Z|Y).
\end{split}
\end{equation*}

\subsection{Extensión}
\label{\detokenize{02_mathematical_foundations/properties_theorems:extension}}
\sphinxAtStartPar
Para múltiples variables:
\begin{equation*}
\begin{split}
I(X;Y_1,\dots,Y_n)=
\sum_{i=1}^{n}
I(X;Y_i|Y_1,\dots,Y_{i-1}).
\end{split}
\end{equation*}

\bigskip\hrule\bigskip



\section{Desigualdad de Procesamiento de Datos}
\label{\detokenize{02_mathematical_foundations/properties_theorems:desigualdad-de-procesamiento-de-datos}}

\subsection{Teorema}
\label{\detokenize{02_mathematical_foundations/properties_theorems:id11}}
\sphinxAtStartPar
Si:
\begin{equation*}
\begin{split}
X\rightarrow Y\rightarrow Z
\end{split}
\end{equation*}
\sphinxAtStartPar
es una cadena de Markov, entonces:
\begin{equation*}
\begin{split}
I(X;Z)\le I(X;Y).
\end{split}
\end{equation*}

\subsection{Interpretación}
\label{\detokenize{02_mathematical_foundations/properties_theorems:id12}}
\sphinxAtStartPar
Ninguna transformación puede aumentar la información.

\sphinxAtStartPar
Aplicaciones:
\begin{itemize}
\item {} 
\sphinxAtStartPar
redes neuronales profundas

\item {} 
\sphinxAtStartPar
compresión de datos

\item {} 
\sphinxAtStartPar
modelos generativos

\end{itemize}


\bigskip\hrule\bigskip



\section{Monotonicidad Condicional}
\label{\detokenize{02_mathematical_foundations/properties_theorems:monotonicidad-condicional}}

\subsection{Teorema}
\label{\detokenize{02_mathematical_foundations/properties_theorems:id13}}\begin{equation*}
\begin{split}
H(X|Y,Z)\le H(X|Y).
\end{split}
\end{equation*}
\sphinxAtStartPar
Más información nunca aumenta la incertidumbre.


\bigskip\hrule\bigskip



\section{Descomposición de Información Mutua}
\label{\detokenize{02_mathematical_foundations/properties_theorems:descomposicion-de-informacion-mutua}}

\subsection{Identidad}
\label{\detokenize{02_mathematical_foundations/properties_theorems:identidad}}\begin{equation*}
\begin{split}
I(X;Y)=H(X)-H(X|Y).
\end{split}
\end{equation*}
\sphinxAtStartPar
También:
\begin{equation*}
\begin{split}
I(X;Y)=H(Y)-H(Y|X).
\end{split}
\end{equation*}

\bigskip\hrule\bigskip



\section{Submodularidad}
\label{\detokenize{02_mathematical_foundations/properties_theorems:submodularidad}}

\subsection{Teorema}
\label{\detokenize{02_mathematical_foundations/properties_theorems:id14}}
\sphinxAtStartPar
La entropía es una función submodular:
\begin{equation*}
\begin{split}
H(A)+H(B)\ge H(A\cup B)+H(A\cap B).
\end{split}
\end{equation*}
\sphinxAtStartPar
Importancia:
\begin{itemize}
\item {} 
\sphinxAtStartPar
selección de variables

\item {} 
\sphinxAtStartPar
optimización combinatoria

\item {} 
\sphinxAtStartPar
aprendizaje estructural

\end{itemize}


\bigskip\hrule\bigskip



\section{Convexidad y Concavidad}
\label{\detokenize{02_mathematical_foundations/properties_theorems:convexidad-y-concavidad}}\begin{itemize}
\item {} 
\sphinxAtStartPar
La entropía es función cóncava de la distribución

\item {} 
\sphinxAtStartPar
La divergencia KL es convexa

\end{itemize}

\sphinxAtStartPar
Consecuencias:
\begin{itemize}
\item {} 
\sphinxAtStartPar
estabilidad numérica

\item {} 
\sphinxAtStartPar
existencia de soluciones óptimas

\item {} 
\sphinxAtStartPar
optimización convexa

\end{itemize}


\bigskip\hrule\bigskip



\section{Información Total (Dependencia Multivariada)}
\label{\detokenize{02_mathematical_foundations/properties_theorems:informacion-total-dependencia-multivariada}}

\subsection{Definición}
\label{\detokenize{02_mathematical_foundations/properties_theorems:definicion}}\begin{equation*}
\begin{split}
TC(X_1,\dots,X_n)=
D_{KL}(P\|\prod_i P_i).
\end{split}
\end{equation*}
\sphinxAtStartPar
Propiedades:
\begin{itemize}
\item {} 
\sphinxAtStartPar
siempre no negativa

\item {} 
\sphinxAtStartPar
cero solo si independencia total

\item {} 
\sphinxAtStartPar
generalización multivariada de información mutua

\end{itemize}


\bigskip\hrule\bigskip



\section{Desigualdades Multivariadas}
\label{\detokenize{02_mathematical_foundations/properties_theorems:desigualdades-multivariadas}}
\sphinxAtStartPar
Para variables múltiples:
\begin{equation*}
\begin{split}
I(X;Y,Z)\ge I(X;Y).
\end{split}
\end{equation*}\begin{equation*}
\begin{split}
H(X,Y,Z)\le H(X)+H(Y)+H(Z).
\end{split}
\end{equation*}

\bigskip\hrule\bigskip



\section{Ejemplo Numérico}
\label{\detokenize{02_mathematical_foundations/properties_theorems:ejemplo-numerico}}
\sphinxAtStartPar
Sea:
\begin{equation*}
\begin{split}
P(0,0)=0.4,\;
P(0,1)=0.1,\;
P(1,0)=0.2,\;
P(1,1)=0.3.
\end{split}
\end{equation*}
\sphinxAtStartPar
Entonces:
\begin{equation*}
\begin{split}
I(X;Y)\ge 0
\end{split}
\end{equation*}
\sphinxAtStartPar
y
\begin{equation*}
\begin{split}
H(X,Y)\le H(X)+H(Y).
\end{split}
\end{equation*}
\sphinxAtStartPar
Este ejemplo se utilizará posteriormente en:
\begin{itemize}
\item {} 
\sphinxAtStartPar
matrices de información mutua

\item {} 
\sphinxAtStartPar
clustering informacional

\item {} 
\sphinxAtStartPar
selección de variables

\end{itemize}


\bigskip\hrule\bigskip



\section{Conclusión}
\label{\detokenize{02_mathematical_foundations/properties_theorems:conclusion}}
\sphinxAtStartPar
Las propiedades presentadas establecen la base teórica necesaria
para desarrollar:
\begin{itemize}
\item {} 
\sphinxAtStartPar
modelos multivariados de dependencia

\item {} 
\sphinxAtStartPar
matrices de información mutua

\item {} 
\sphinxAtStartPar
aprendizaje profundo informacional

\item {} 
\sphinxAtStartPar
algoritmos de ciencia de datos basados en teoría de la información

\end{itemize}

\sphinxAtStartPar
El siguiente capítulo introducirá formalmente la divergencia de
Kullback–Leibler como herramienta central para medir diferencias
entre distribuciones.

\sphinxAtStartPar
{[}\hyperlink{cite.05_appendices/references:id10}{Csiszár, 2011}{]}, {[}\hyperlink{cite.05_appendices/references:id2}{Cover and Thomas, 2006}{]}

\sphinxstepscope


\chapter{Divergencia de Kullback–Leibler}
\label{\detokenize{02_mathematical_foundations/kl_divergence:divergencia-de-kullbackleibler}}\label{\detokenize{02_mathematical_foundations/kl_divergence::doc}}

\section{Introducción}
\label{\detokenize{02_mathematical_foundations/kl_divergence:introduccion}}
\sphinxAtStartPar
La divergencia de Kullback–Leibler (KL) es una medida fundamental
de la diferencia entre dos distribuciones de probabilidad.

\sphinxAtStartPar
En teoría de la información, cuantifica la cantidad de información
perdida al aproximar una distribución verdadera mediante otra.

\sphinxAtStartPar
En ciencia de datos y aprendizaje automático aparece en:
\begin{itemize}
\item {} 
\sphinxAtStartPar
aprendizaje profundo

\item {} 
\sphinxAtStartPar
modelos generativos

\item {} 
\sphinxAtStartPar
inferencia variacional

\item {} 
\sphinxAtStartPar
selección de modelos

\item {} 
\sphinxAtStartPar
regularización informacional

\end{itemize}


\bigskip\hrule\bigskip



\section{Definición Formal (Variables Discretas)}
\label{\detokenize{02_mathematical_foundations/kl_divergence:definicion-formal-variables-discretas}}
\sphinxAtStartPar
Sean \(P\) y \(Q\) dos distribuciones sobre el mismo espacio.

\sphinxAtStartPar
La divergencia KL se define como:
\begin{equation*}
\begin{split}
D_{KL}(P\|Q)=
\sum_x P(x)\log\frac{P(x)}{Q(x)}.
\end{split}
\end{equation*}
\sphinxAtStartPar
Interpretación:
\begin{itemize}
\item {} 
\sphinxAtStartPar
mide discrepancia probabilística

\item {} 
\sphinxAtStartPar
no es una distancia métrica

\item {} 
\sphinxAtStartPar
es una medida direccional

\end{itemize}


\bigskip\hrule\bigskip



\section{Definición Continua}
\label{\detokenize{02_mathematical_foundations/kl_divergence:definicion-continua}}
\sphinxAtStartPar
Para variables continuas:
\begin{equation*}
\begin{split}
D_{KL}(P\|Q)=
\int p(x)\log\frac{p(x)}{q(x)}dx.
\end{split}
\end{equation*}

\bigskip\hrule\bigskip



\section{Interpretación Informacional}
\label{\detokenize{02_mathematical_foundations/kl_divergence:interpretacion-informacional}}
\sphinxAtStartPar
Reescribiendo:
\begin{equation*}
\begin{split}
D_{KL}(P\|Q)
=
-\sum P(x)\log Q(x)
-
H(P).
\end{split}
\end{equation*}
\sphinxAtStartPar
Interpretación:
\begin{itemize}
\item {} 
\sphinxAtStartPar
diferencia entre entropía cruzada y entropía real

\item {} 
\sphinxAtStartPar
exceso de longitud de código

\item {} 
\sphinxAtStartPar
pérdida de eficiencia de compresión

\end{itemize}


\bigskip\hrule\bigskip



\section{Propiedades Fundamentales}
\label{\detokenize{02_mathematical_foundations/kl_divergence:propiedades-fundamentales}}

\subsection{No Negatividad}
\label{\detokenize{02_mathematical_foundations/kl_divergence:no-negatividad}}\begin{equation*}
\begin{split}
D_{KL}(P\|Q)\ge 0.
\end{split}
\end{equation*}
\sphinxAtStartPar
Igualdad solo si:
\begin{equation*}
\begin{split}
P=Q.
\end{split}
\end{equation*}

\bigskip\hrule\bigskip



\subsection{No Simetría}
\label{\detokenize{02_mathematical_foundations/kl_divergence:no-simetria}}\begin{equation*}
\begin{split}
D_{KL}(P\|Q)\neq D_{KL}(Q\|P).
\end{split}
\end{equation*}
\sphinxAtStartPar
Consecuencia:

\sphinxAtStartPar
no es una métrica.


\bigskip\hrule\bigskip



\subsection{No Cumple Desigualdad Triangular}
\label{\detokenize{02_mathematical_foundations/kl_divergence:no-cumple-desigualdad-triangular}}
\sphinxAtStartPar
No define distancia geométrica clásica.


\bigskip\hrule\bigskip



\subsection{Convexidad}
\label{\detokenize{02_mathematical_foundations/kl_divergence:convexidad}}
\sphinxAtStartPar
La divergencia KL es convexa respecto a \(Q\).

\sphinxAtStartPar
Importancia:
\begin{itemize}
\item {} 
\sphinxAtStartPar
optimización convexa

\item {} 
\sphinxAtStartPar
aprendizaje probabilístico

\item {} 
\sphinxAtStartPar
estabilidad numérica

\end{itemize}


\bigskip\hrule\bigskip



\section{Relación con Entropía Cruzada}
\label{\detokenize{02_mathematical_foundations/kl_divergence:relacion-con-entropia-cruzada}}
\sphinxAtStartPar
Definición:
\begin{equation*}
\begin{split}
H(P,Q)=
-\sum P(x)\log Q(x).
\end{split}
\end{equation*}
\sphinxAtStartPar
Entonces:
\begin{equation*}
\begin{split}
D_{KL}(P\|Q)=H(P,Q)-H(P).
\end{split}
\end{equation*}

\bigskip\hrule\bigskip



\section{Relación con Información Mutua}
\label{\detokenize{02_mathematical_foundations/kl_divergence:relacion-con-informacion-mutua}}
\sphinxAtStartPar
Recordemos:
\begin{equation*}
\begin{split}
I(X;Y)=
D_{KL}(P_{X,Y}\|P_XP_Y).
\end{split}
\end{equation*}
\sphinxAtStartPar
Interpretación:

\sphinxAtStartPar
la información mutua mide desviación respecto a independencia.


\bigskip\hrule\bigskip



\section{Ejemplo Discreto}
\label{\detokenize{02_mathematical_foundations/kl_divergence:ejemplo-discreto}}
\sphinxAtStartPar
Sea:
\begin{equation*}
\begin{split}
P=(0.6,0.4),
\quad
Q=(0.5,0.5).
\end{split}
\end{equation*}
\sphinxAtStartPar
Entonces:
\begin{equation*}
\begin{split}
D_{KL}(P\|Q)
=
0.6\log\frac{0.6}{0.5}
+
0.4\log\frac{0.4}{0.5}.
\end{split}
\end{equation*}
\sphinxAtStartPar
Calculando:
\begin{equation*}
\begin{split}
=0.6\log(1.2)+0.4\log(0.8).
\end{split}
\end{equation*}
\sphinxAtStartPar
Resultado positivo ⇒ distribuciones diferentes.


\bigskip\hrule\bigskip



\section{Ejemplo Interpretativo (Modelo Mal Especificado)}
\label{\detokenize{02_mathematical_foundations/kl_divergence:ejemplo-interpretativo-modelo-mal-especificado}}
\sphinxAtStartPar
Supongamos:
\begin{itemize}
\item {} 
\sphinxAtStartPar
distribución real: \(P\)

\item {} 
\sphinxAtStartPar
modelo estimado: \(Q\)

\end{itemize}

\sphinxAtStartPar
La divergencia KL mide:
\begin{itemize}
\item {} 
\sphinxAtStartPar
error informacional del modelo

\item {} 
\sphinxAtStartPar
pérdida de eficiencia predictiva

\end{itemize}

\sphinxAtStartPar
Aplicaciones:
\begin{itemize}
\item {} 
\sphinxAtStartPar
modelos bayesianos

\item {} 
\sphinxAtStartPar
machine learning

\item {} 
\sphinxAtStartPar
selección de arquitectura

\end{itemize}


\bigskip\hrule\bigskip



\section{Divergencia KL Condicional}
\label{\detokenize{02_mathematical_foundations/kl_divergence:divergencia-kl-condicional}}
\sphinxAtStartPar
Para variables:
\begin{equation*}
\begin{split}
D_{KL}(P(X|Y)\|Q(X|Y)).
\end{split}
\end{equation*}
\sphinxAtStartPar
Definición promedio:
\begin{equation*}
\begin{split}
\mathbb{E}_Y
\left[
D_{KL}(P(X|Y)\|Q(X|Y))
\right].
\end{split}
\end{equation*}
\sphinxAtStartPar
Uso:
\begin{itemize}
\item {} 
\sphinxAtStartPar
modelos jerárquicos

\item {} 
\sphinxAtStartPar
inferencia variacional

\item {} 
\sphinxAtStartPar
redes probabilísticas

\end{itemize}


\bigskip\hrule\bigskip



\section{Descomposición de KL}
\label{\detokenize{02_mathematical_foundations/kl_divergence:descomposicion-de-kl}}
\sphinxAtStartPar
Para variables conjuntas:
\begin{equation*}
\begin{split}
D_{KL}(P(X,Y)\|Q(X,Y))
=
D_{KL}(P(X)\|Q(X))
+
D_{KL}(P(Y|X)\|Q(Y|X)).
\end{split}
\end{equation*}
\sphinxAtStartPar
Importancia:
\begin{itemize}
\item {} 
\sphinxAtStartPar
modelos estructurados

\item {} 
\sphinxAtStartPar
factorización probabilística

\item {} 
\sphinxAtStartPar
aprendizaje profundo

\end{itemize}


\bigskip\hrule\bigskip



\section{KL en Aprendizaje Automático}
\label{\detokenize{02_mathematical_foundations/kl_divergence:kl-en-aprendizaje-automatico}}

\subsection{Función de Pérdida}
\label{\detokenize{02_mathematical_foundations/kl_divergence:funcion-de-perdida}}
\sphinxAtStartPar
Muchos modelos minimizan:
\begin{equation*}
\begin{split}
D_{KL}(P_{datos}\|P_{modelo}).
\end{split}
\end{equation*}
\sphinxAtStartPar
Ejemplos:
\begin{itemize}
\item {} 
\sphinxAtStartPar
modelos generativos

\item {} 
\sphinxAtStartPar
VAE

\item {} 
\sphinxAtStartPar
language models

\end{itemize}


\bigskip\hrule\bigskip



\subsection{Regularización}
\label{\detokenize{02_mathematical_foundations/kl_divergence:regularizacion}}
\sphinxAtStartPar
KL controla:
\begin{itemize}
\item {} 
\sphinxAtStartPar
complejidad del modelo

\item {} 
\sphinxAtStartPar
sobreajuste

\item {} 
\sphinxAtStartPar
divergencia respecto a prior

\end{itemize}


\bigskip\hrule\bigskip



\section{Inferencia Variacional}
\label{\detokenize{02_mathematical_foundations/kl_divergence:inferencia-variacional}}
\sphinxAtStartPar
Objetivo:

\sphinxAtStartPar
aproximar posterior:
\begin{equation*}
\begin{split}
p(z|x)
\end{split}
\end{equation*}
\sphinxAtStartPar
mediante:
\begin{equation*}
\begin{split}
q(z).
\end{split}
\end{equation*}
\sphinxAtStartPar
Minimizando:
\begin{equation*}
\begin{split}
D_{KL}(q(z)\|p(z|x)).
\end{split}
\end{equation*}
\sphinxAtStartPar
Base matemática de:
\begin{itemize}
\item {} 
\sphinxAtStartPar
autoencoders variacionales

\item {} 
\sphinxAtStartPar
modelos bayesianos modernos

\end{itemize}


\bigskip\hrule\bigskip



\section{Relación con Otras Medidas}
\label{\detokenize{02_mathematical_foundations/kl_divergence:relacion-con-otras-medidas}}\begin{itemize}
\item {} 
\sphinxAtStartPar
Jensen–Shannon divergence

\item {} 
\sphinxAtStartPar
Total variation

\item {} 
\sphinxAtStartPar
Wasserstein distance

\end{itemize}

\sphinxAtStartPar
KL destaca por:
\begin{itemize}
\item {} 
\sphinxAtStartPar
tractabilidad analítica

\item {} 
\sphinxAtStartPar
diferenciabilidad

\item {} 
\sphinxAtStartPar
conexión directa con entropía

\end{itemize}


\bigskip\hrule\bigskip



\section{Ejemplo Multivariado}
\label{\detokenize{02_mathematical_foundations/kl_divergence:ejemplo-multivariado}}
\sphinxAtStartPar
Sea:
\begin{equation*}
\begin{split}
P(X,Y)
\end{split}
\end{equation*}
\sphinxAtStartPar
y modelo independiente:
\begin{equation*}
\begin{split}
Q=P_XP_Y.
\end{split}
\end{equation*}
\sphinxAtStartPar
Entonces:
\begin{equation*}
\begin{split}
D_{KL}(P\|Q)=I(X;Y).
\end{split}
\end{equation*}
\sphinxAtStartPar
Interpretación:

\sphinxAtStartPar
la información mutua es KL contra independencia.


\bigskip\hrule\bigskip



\section{Conclusión}
\label{\detokenize{02_mathematical_foundations/kl_divergence:conclusion}}
\sphinxAtStartPar
La divergencia de Kullback–Leibler constituye el núcleo matemático
de la teoría moderna de la información y el aprendizaje automático.

\sphinxAtStartPar
Proporciona:
\begin{itemize}
\item {} 
\sphinxAtStartPar
medida fundamental de discrepancia

\item {} 
\sphinxAtStartPar
base para información mutua

\item {} 
\sphinxAtStartPar
herramienta central en deep learning

\item {} 
\sphinxAtStartPar
fundamento de inferencia probabilística

\end{itemize}

\sphinxAtStartPar
En los siguientes capítulos se utilizará extensivamente para:
\begin{itemize}
\item {} 
\sphinxAtStartPar
matrices de información mutua

\item {} 
\sphinxAtStartPar
modelos multivariados

\item {} 
\sphinxAtStartPar
selección de características

\item {} 
\sphinxAtStartPar
análisis informacional en ciencia de datos

\end{itemize}

\sphinxAtStartPar
{[}\hyperlink{cite.05_appendices/references:id15}{Kullback and Leibler, 1951}{]}, {[}\hyperlink{cite.05_appendices/references:id4}{Bishop, 2006}{]}

\sphinxstepscope


\part{Modelos Multivariados}

\sphinxstepscope


\chapter{Matriz de Información Mutua}
\label{\detokenize{03_multivariate_models/mi_matrix:matriz-de-informacion-mutua}}\label{\detokenize{03_multivariate_models/mi_matrix::doc}}

\section{Introducción}
\label{\detokenize{03_multivariate_models/mi_matrix:introduccion}}
\sphinxAtStartPar
En sistemas multivariados resulta fundamental cuantificar
las dependencias entre múltiples variables aleatorias.

\sphinxAtStartPar
Mientras que la información mutua clásica mide la dependencia
entre dos variables, la matriz de información mutua permite
extender este concepto a conjuntos de variables.

\sphinxAtStartPar
Aplicaciones principales:
\begin{itemize}
\item {} 
\sphinxAtStartPar
análisis multivariado

\item {} 
\sphinxAtStartPar
selección de características

\item {} 
\sphinxAtStartPar
aprendizaje automático

\item {} 
\sphinxAtStartPar
análisis de redes complejas

\item {} 
\sphinxAtStartPar
modelamiento probabilístico de alta dimensión

\end{itemize}


\bigskip\hrule\bigskip



\section{Definición del Sistema Multivariado}
\label{\detokenize{03_multivariate_models/mi_matrix:definicion-del-sistema-multivariado}}
\sphinxAtStartPar
Sea un vector aleatorio:
\begin{equation*}
\begin{split}
\mathbf{X} = (X_1,X_2,\dots,X_n).
\end{split}
\end{equation*}
\sphinxAtStartPar
Cada variable tiene distribución marginal:
\begin{equation*}
\begin{split}
P_{X_i}.
\end{split}
\end{equation*}
\sphinxAtStartPar
La distribución conjunta es:
\begin{equation*}
\begin{split}
P_{X_1,\dots,X_n}.
\end{split}
\end{equation*}

\bigskip\hrule\bigskip



\section{Definición de la Matriz de Información Mutua}
\label{\detokenize{03_multivariate_models/mi_matrix:definicion-de-la-matriz-de-informacion-mutua}}
\sphinxAtStartPar
La matriz de información mutua se define como:
\begin{equation*}
\begin{split}
\mathbf{M} =
[I(X_i;X_j)]_{i,j=1}^n.
\end{split}
\end{equation*}
\sphinxAtStartPar
Donde cada entrada es:
\begin{equation*}
\begin{split}
M_{ij}=I(X_i;X_j).
\end{split}
\end{equation*}

\bigskip\hrule\bigskip



\section{Estructura de la Matriz}
\label{\detokenize{03_multivariate_models/mi_matrix:estructura-de-la-matriz}}

\subsection{Diagonal}
\label{\detokenize{03_multivariate_models/mi_matrix:diagonal}}\begin{equation*}
\begin{split}
M_{ii}=H(X_i).
\end{split}
\end{equation*}
\sphinxAtStartPar
Interpretación:

\sphinxAtStartPar
la información mutua de una variable consigo misma
equivale a su entropía.


\bigskip\hrule\bigskip



\subsection{Simetría}
\label{\detokenize{03_multivariate_models/mi_matrix:simetria}}\begin{equation*}
\begin{split}
I(X_i;X_j)=I(X_j;X_i).
\end{split}
\end{equation*}
\sphinxAtStartPar
Por lo tanto:
\begin{equation*}
\begin{split}
\mathbf{M}
\end{split}
\end{equation*}
\sphinxAtStartPar
es una matriz simétrica.


\bigskip\hrule\bigskip



\section{Interpretación Informacional}
\label{\detokenize{03_multivariate_models/mi_matrix:interpretacion-informacional}}
\sphinxAtStartPar
Cada elemento mide:
\begin{itemize}
\item {} 
\sphinxAtStartPar
reducción de incertidumbre

\item {} 
\sphinxAtStartPar
dependencia estadística

\item {} 
\sphinxAtStartPar
redundancia informacional

\item {} 
\sphinxAtStartPar
relación funcional potencial

\end{itemize}

\sphinxAtStartPar
Valores altos ⇒ fuerte dependencia.

\sphinxAtStartPar
Valores cercanos a cero ⇒ independencia.


\bigskip\hrule\bigskip



\section{Relación con la Matriz de Covarianza}
\label{\detokenize{03_multivariate_models/mi_matrix:relacion-con-la-matriz-de-covarianza}}

\begin{savenotes}\sphinxattablestart
\centering
\begin{tabulary}{\linewidth}[t]{|T|T|T|}
\hline
\sphinxstyletheadfamily 
\sphinxAtStartPar
Aspecto
&\sphinxstyletheadfamily 
\sphinxAtStartPar
Covarianza
&\sphinxstyletheadfamily 
\sphinxAtStartPar
Información Mutua
\\
\hline
\sphinxAtStartPar
Tipo de dependencia
&
\sphinxAtStartPar
Lineal
&
\sphinxAtStartPar
General
\\
\hline
\sphinxAtStartPar
Sensibilidad
&
\sphinxAtStartPar
Baja a no linealidad
&
\sphinxAtStartPar
Alta
\\
\hline
\sphinxAtStartPar
Invarianza a transformaciones
&
\sphinxAtStartPar
Limitada
&
\sphinxAtStartPar
Más robusta
\\
\hline
\sphinxAtStartPar
Aplicación
&
\sphinxAtStartPar
Estadística clásica
&
\sphinxAtStartPar
Ciencia de datos moderna
\\
\hline
\end{tabulary}
\par
\sphinxattableend\end{savenotes}

\sphinxAtStartPar
La matriz de información mutua detecta:
\begin{itemize}
\item {} 
\sphinxAtStartPar
dependencias no lineales

\item {} 
\sphinxAtStartPar
relaciones complejas

\item {} 
\sphinxAtStartPar
interacciones ocultas

\end{itemize}


\bigskip\hrule\bigskip



\section{Construcción Práctica}
\label{\detokenize{03_multivariate_models/mi_matrix:construccion-practica}}

\subsection{Paso 1: Estimar distribuciones}
\label{\detokenize{03_multivariate_models/mi_matrix:paso-1-estimar-distribuciones}}\begin{itemize}
\item {} 
\sphinxAtStartPar
histogramas

\item {} 
\sphinxAtStartPar
KDE

\item {} 
\sphinxAtStartPar
estimadores k\sphinxhyphen{}NN

\item {} 
\sphinxAtStartPar
modelos probabilísticos

\end{itemize}


\subsection{Paso 2: Calcular información mutua}
\label{\detokenize{03_multivariate_models/mi_matrix:paso-2-calcular-informacion-mutua}}\begin{equation*}
\begin{split}
I(X_i;X_j)
=
H(X_i)+H(X_j)-H(X_i,X_j).
\end{split}
\end{equation*}

\subsection{Paso 3: Construir la matriz}
\label{\detokenize{03_multivariate_models/mi_matrix:paso-3-construir-la-matriz}}
\sphinxAtStartPar
Para todo par:
\begin{equation*}
\begin{split}
(i,j).
\end{split}
\end{equation*}

\bigskip\hrule\bigskip



\section{Ejemplo Conceptual}
\label{\detokenize{03_multivariate_models/mi_matrix:ejemplo-conceptual}}
\sphinxAtStartPar
Sea:
\begin{equation*}
\begin{split}
(X_1,X_2,X_3).
\end{split}
\end{equation*}
\sphinxAtStartPar
La matriz es:
\begin{equation*}
\begin{split}
\begin{pmatrix}
H(X_1) & I(X_1;X_2) & I(X_1;X_3) \\
I(X_2;X_1) & H(X_2) & I(X_2;X_3) \\
I(X_3;X_1) & I(X_3;X_2) & H(X_3)
\end{pmatrix}.
\end{split}
\end{equation*}

\bigskip\hrule\bigskip



\section{Propiedades Matemáticas}
\label{\detokenize{03_multivariate_models/mi_matrix:propiedades-matematicas}}

\subsection{No Negatividad}
\label{\detokenize{03_multivariate_models/mi_matrix:no-negatividad}}\begin{equation*}
\begin{split}
M_{ij}\ge 0.
\end{split}
\end{equation*}

\bigskip\hrule\bigskip



\subsection{Simetría}
\label{\detokenize{03_multivariate_models/mi_matrix:id1}}\begin{equation*}
\begin{split}
\mathbf{M}=\mathbf{M}^T.
\end{split}
\end{equation*}

\bigskip\hrule\bigskip



\subsection{Dependencia Funcional}
\label{\detokenize{03_multivariate_models/mi_matrix:dependencia-funcional}}
\sphinxAtStartPar
Si:
\begin{equation*}
\begin{split}
X_j=f(X_i)
\end{split}
\end{equation*}
\sphinxAtStartPar
entonces:
\begin{equation*}
\begin{split}
I(X_i;X_j)=H(X_j).
\end{split}
\end{equation*}

\bigskip\hrule\bigskip



\section{Interpretación Geométrica}
\label{\detokenize{03_multivariate_models/mi_matrix:interpretacion-geometrica}}
\sphinxAtStartPar
La matriz define un grafo informacional:
\begin{itemize}
\item {} 
\sphinxAtStartPar
nodos = variables

\item {} 
\sphinxAtStartPar
pesos = información mutua

\end{itemize}

\sphinxAtStartPar
Esto permite:
\begin{itemize}
\item {} 
\sphinxAtStartPar
análisis de redes

\item {} 
\sphinxAtStartPar
clustering estructural

\item {} 
\sphinxAtStartPar
inferencia causal preliminar

\end{itemize}


\bigskip\hrule\bigskip



\section{Uso en Ciencia de Datos}
\label{\detokenize{03_multivariate_models/mi_matrix:uso-en-ciencia-de-datos}}

\subsection{Selección de Variables}
\label{\detokenize{03_multivariate_models/mi_matrix:seleccion-de-variables}}
\sphinxAtStartPar
Identificar:
\begin{itemize}
\item {} 
\sphinxAtStartPar
redundancia

\item {} 
\sphinxAtStartPar
variables irrelevantes

\item {} 
\sphinxAtStartPar
grupos altamente dependientes

\end{itemize}


\bigskip\hrule\bigskip



\subsection{Modelos Generativos}
\label{\detokenize{03_multivariate_models/mi_matrix:modelos-generativos}}
\sphinxAtStartPar
Detectar:
\begin{itemize}
\item {} 
\sphinxAtStartPar
estructuras probabilísticas

\item {} 
\sphinxAtStartPar
correlaciones ocultas

\item {} 
\sphinxAtStartPar
dependencias jerárquicas

\end{itemize}


\bigskip\hrule\bigskip



\subsection{Deep Learning}
\label{\detokenize{03_multivariate_models/mi_matrix:deep-learning}}\begin{itemize}
\item {} 
\sphinxAtStartPar
análisis de representaciones internas

\item {} 
\sphinxAtStartPar
aprendizaje auto\sphinxhyphen{}supervisado

\item {} 
\sphinxAtStartPar
regularización informacional

\end{itemize}


\bigskip\hrule\bigskip



\section{Ejemplo Conceptual de Ciencia de Datos}
\label{\detokenize{03_multivariate_models/mi_matrix:ejemplo-conceptual-de-ciencia-de-datos}}
\sphinxAtStartPar
Variables:
\begin{itemize}
\item {} 
\sphinxAtStartPar
Edad

\item {} 
\sphinxAtStartPar
Presión arterial

\item {} 
\sphinxAtStartPar
Índice de masa corporal

\item {} 
\sphinxAtStartPar
Nivel de actividad

\end{itemize}

\sphinxAtStartPar
La matriz permite:
\begin{itemize}
\item {} 
\sphinxAtStartPar
detectar redundancia clínica

\item {} 
\sphinxAtStartPar
analizar interacción fisiológica

\item {} 
\sphinxAtStartPar
reducir dimensionalidad

\end{itemize}


\bigskip\hrule\bigskip



\section{Limitaciones}
\label{\detokenize{03_multivariate_models/mi_matrix:limitaciones}}\begin{itemize}
\item {} 
\sphinxAtStartPar
estimación difícil en alta dimensión

\item {} 
\sphinxAtStartPar
sesgo en muestras pequeñas

\item {} 
\sphinxAtStartPar
necesidad de estimadores robustos

\item {} 
\sphinxAtStartPar
costo computacional elevado

\end{itemize}


\bigskip\hrule\bigskip



\section{Extensiones}
\label{\detokenize{03_multivariate_models/mi_matrix:extensiones}}\begin{itemize}
\item {} 
\sphinxAtStartPar
matriz de información mutua condicional

\item {} 
\sphinxAtStartPar
grafos de información

\item {} 
\sphinxAtStartPar
redes bayesianas

\item {} 
\sphinxAtStartPar
matrices normalizadas

\end{itemize}


\bigskip\hrule\bigskip



\section{Conclusión}
\label{\detokenize{03_multivariate_models/mi_matrix:conclusion}}
\sphinxAtStartPar
La matriz de información mutua generaliza el concepto
de dependencia informacional al contexto multivariado.

\sphinxAtStartPar
Proporciona una herramienta poderosa para:
\begin{itemize}
\item {} 
\sphinxAtStartPar
análisis estadístico avanzado

\item {} 
\sphinxAtStartPar
aprendizaje automático

\item {} 
\sphinxAtStartPar
ciencia de datos moderna

\item {} 
\sphinxAtStartPar
modelamiento probabilístico de alta dimensión

\end{itemize}

\sphinxAtStartPar
En el siguiente capítulo se desarrollará la
información mutua condicional multivariada,
clave para analizar dependencias indirectas
y relaciones causales complejas.

\sphinxAtStartPar
{[}\hyperlink{cite.05_appendices/references:id5}{Murphy, 2012}{]}

\sphinxstepscope


\chapter{Información Mutua Condicional Multivariada}
\label{\detokenize{03_multivariate_models/conditional_mi:informacion-mutua-condicional-multivariada}}\label{\detokenize{03_multivariate_models/conditional_mi::doc}}

\section{Introducción}
\label{\detokenize{03_multivariate_models/conditional_mi:introduccion}}
\sphinxAtStartPar
En sistemas multivariados complejos, las dependencias entre variables
no siempre son directas. Muchas relaciones observadas pueden deberse
a variables intermedias o factores latentes.

\sphinxAtStartPar
La información mutua condicional permite medir la dependencia entre
dos variables eliminando la influencia de una o más variables de control.

\sphinxAtStartPar
Este concepto es fundamental en:
\begin{itemize}
\item {} 
\sphinxAtStartPar
inferencia causal

\item {} 
\sphinxAtStartPar
redes bayesianas

\item {} 
\sphinxAtStartPar
selección de variables

\item {} 
\sphinxAtStartPar
análisis de dependencia estructural

\item {} 
\sphinxAtStartPar
aprendizaje automático moderno

\end{itemize}


\bigskip\hrule\bigskip



\section{Definición General}
\label{\detokenize{03_multivariate_models/conditional_mi:definicion-general}}
\sphinxAtStartPar
Sean variables aleatorias:
\begin{equation*}
\begin{split}
X,\ Y,\ Z.
\end{split}
\end{equation*}
\sphinxAtStartPar
La información mutua condicional se define como:
\begin{equation*}
\begin{split}
I(X;Y|Z)
=
H(X|Z)
-
H(X|Y,Z).
\end{split}
\end{equation*}
\sphinxAtStartPar
También puede escribirse como:
\begin{equation*}
\begin{split}
I(X;Y|Z)
=
H(X,Z)
+
H(Y,Z)
-
H(Z)
-
H(X,Y,Z).
\end{split}
\end{equation*}

\bigskip\hrule\bigskip



\section{Interpretación Informacional}
\label{\detokenize{03_multivariate_models/conditional_mi:interpretacion-informacional}}
\sphinxAtStartPar
La cantidad:
\begin{equation*}
\begin{split}
I(X;Y|Z)
\end{split}
\end{equation*}
\sphinxAtStartPar
mide la reducción de incertidumbre sobre \(X\) al conocer \(Y\),
una vez que ya conocemos \(Z\).

\sphinxAtStartPar
Interpretación:
\begin{itemize}
\item {} 
\sphinxAtStartPar
dependencia directa residual

\item {} 
\sphinxAtStartPar
información adicional

\item {} 
\sphinxAtStartPar
asociación libre de confusores

\end{itemize}


\bigskip\hrule\bigskip



\section{Caso Multivariado General}
\label{\detokenize{03_multivariate_models/conditional_mi:caso-multivariado-general}}
\sphinxAtStartPar
Sean vectores aleatorios:
\begin{equation*}
\begin{split}
\mathbf{X},\ \mathbf{Y},\ \mathbf{Z}.
\end{split}
\end{equation*}
\sphinxAtStartPar
La definición general es:
\begin{equation*}
\begin{split}
I(\mathbf{X};\mathbf{Y}|\mathbf{Z})
=
H(\mathbf{X}|\mathbf{Z})
-
H(\mathbf{X}|\mathbf{Y},\mathbf{Z}).
\end{split}
\end{equation*}
\sphinxAtStartPar
Equivalentemente:
\begin{equation*}
\begin{split}
I(\mathbf{X};\mathbf{Y}|\mathbf{Z})
=
H(\mathbf{X},\mathbf{Z})
+
H(\mathbf{Y},\mathbf{Z})
-
H(\mathbf{Z})
-
H(\mathbf{X},\mathbf{Y},\mathbf{Z}).
\end{split}
\end{equation*}

\bigskip\hrule\bigskip



\section{Interpretación Geométrica}
\label{\detokenize{03_multivariate_models/conditional_mi:interpretacion-geometrica}}
\sphinxAtStartPar
En términos de incertidumbre:
\begin{itemize}
\item {} 
\sphinxAtStartPar
\(Z\) define el contexto informacional

\item {} 
\sphinxAtStartPar
\(Y\) aporta información adicional sobre \(X\)

\item {} 
\sphinxAtStartPar
la dependencia medida es la parte irreducible
una vez controlado \(Z\)

\end{itemize}


\bigskip\hrule\bigskip



\section{Independencia Condicional}
\label{\detokenize{03_multivariate_models/conditional_mi:independencia-condicional}}
\sphinxAtStartPar
Si:
\begin{equation*}
\begin{split}
X \perp Y \mid Z
\end{split}
\end{equation*}
\sphinxAtStartPar
entonces:
\begin{equation*}
\begin{split}
I(X;Y|Z)=0.
\end{split}
\end{equation*}
\sphinxAtStartPar
Esto permite detectar:
\begin{itemize}
\item {} 
\sphinxAtStartPar
independencia estructural

\item {} 
\sphinxAtStartPar
separación d\sphinxhyphen{}separation en grafos

\item {} 
\sphinxAtStartPar
relaciones indirectas

\end{itemize}


\bigskip\hrule\bigskip



\section{Propiedades Fundamentales}
\label{\detokenize{03_multivariate_models/conditional_mi:propiedades-fundamentales}}

\subsection{No Negatividad}
\label{\detokenize{03_multivariate_models/conditional_mi:no-negatividad}}\begin{equation*}
\begin{split}
I(X;Y|Z)\ge 0.
\end{split}
\end{equation*}

\bigskip\hrule\bigskip



\subsection{Simetría}
\label{\detokenize{03_multivariate_models/conditional_mi:simetria}}\begin{equation*}
\begin{split}
I(X;Y|Z)=I(Y;X|Z).
\end{split}
\end{equation*}

\bigskip\hrule\bigskip



\subsection{Regla de Cadena}
\label{\detokenize{03_multivariate_models/conditional_mi:regla-de-cadena}}\begin{equation*}
\begin{split}
I(X;Y,Z)
=
I(X;Z)
+
I(X;Y|Z).
\end{split}
\end{equation*}

\bigskip\hrule\bigskip



\subsection{Expansión Multivariable}
\label{\detokenize{03_multivariate_models/conditional_mi:expansion-multivariable}}
\sphinxAtStartPar
Para variables múltiples:
\begin{equation*}
\begin{split}
I(X;Y_1,Y_2|Z)
=
I(X;Y_1|Z)
+
I(X;Y_2|Y_1,Z).
\end{split}
\end{equation*}

\bigskip\hrule\bigskip



\section{Relación con Grafos Probabilísticos}
\label{\detokenize{03_multivariate_models/conditional_mi:relacion-con-grafos-probabilisticos}}
\sphinxAtStartPar
La información mutua condicional permite:
\begin{itemize}
\item {} 
\sphinxAtStartPar
aprender estructuras de redes bayesianas

\item {} 
\sphinxAtStartPar
identificar aristas redundantes

\item {} 
\sphinxAtStartPar
evaluar independencia estructural

\end{itemize}

\sphinxAtStartPar
Ejemplo:
\begin{equation*}
\begin{split}
X \rightarrow Z \rightarrow Y
\end{split}
\end{equation*}
\sphinxAtStartPar
puede implicar:
\begin{equation*}
\begin{split}
I(X;Y)>0
\end{split}
\end{equation*}
\sphinxAtStartPar
pero:
\begin{equation*}
\begin{split}
I(X;Y|Z)=0.
\end{split}
\end{equation*}

\bigskip\hrule\bigskip



\section{Construcción de la Matriz de Información Mutua Condicional}
\label{\detokenize{03_multivariate_models/conditional_mi:construccion-de-la-matriz-de-informacion-mutua-condicional}}
\sphinxAtStartPar
Sea:
\begin{equation*}
\begin{split}
\mathbf{X}=(X_1,\dots,X_n).
\end{split}
\end{equation*}
\sphinxAtStartPar
Definimos la matriz condicionada a \(Z\) como:
\begin{equation*}
\begin{split}
M_{ij}^{(Z)}=I(X_i;X_j|Z).
\end{split}
\end{equation*}
\sphinxAtStartPar
Aplicaciones:
\begin{itemize}
\item {} 
\sphinxAtStartPar
análisis estructural

\item {} 
\sphinxAtStartPar
inferencia causal

\item {} 
\sphinxAtStartPar
reducción de dimensionalidad

\end{itemize}


\bigskip\hrule\bigskip



\section{Ejemplo Conceptual}
\label{\detokenize{03_multivariate_models/conditional_mi:ejemplo-conceptual}}
\sphinxAtStartPar
Variables:
\begin{itemize}
\item {} 
\sphinxAtStartPar
Edad

\item {} 
\sphinxAtStartPar
Presión arterial

\item {} 
\sphinxAtStartPar
Índice de masa corporal

\item {} 
\sphinxAtStartPar
Actividad física

\end{itemize}

\sphinxAtStartPar
Supongamos:

\sphinxAtStartPar
Edad → Presión
Edad → IMC
IMC → Presión

\sphinxAtStartPar
Entonces:
\begin{equation*}
\begin{split}
I(IMC;Presión)>0
\end{split}
\end{equation*}
\sphinxAtStartPar
pero posiblemente:
\begin{equation*}
\begin{split}
I(IMC;Presión|Edad)
\end{split}
\end{equation*}
\sphinxAtStartPar
disminuye significativamente.


\bigskip\hrule\bigskip



\section{Aplicaciones en Ciencia de Datos}
\label{\detokenize{03_multivariate_models/conditional_mi:aplicaciones-en-ciencia-de-datos}}

\subsection{Selección de Características}
\label{\detokenize{03_multivariate_models/conditional_mi:seleccion-de-caracteristicas}}
\sphinxAtStartPar
Detectar:
\begin{itemize}
\item {} 
\sphinxAtStartPar
redundancia indirecta

\item {} 
\sphinxAtStartPar
variables irrelevantes

\item {} 
\sphinxAtStartPar
efectos mediadores

\end{itemize}


\bigskip\hrule\bigskip



\subsection{Aprendizaje Estructural}
\label{\detokenize{03_multivariate_models/conditional_mi:aprendizaje-estructural}}\begin{itemize}
\item {} 
\sphinxAtStartPar
redes bayesianas

\item {} 
\sphinxAtStartPar
grafos probabilísticos

\item {} 
\sphinxAtStartPar
modelos causales

\end{itemize}


\bigskip\hrule\bigskip



\subsection{Deep Learning}
\label{\detokenize{03_multivariate_models/conditional_mi:deep-learning}}\begin{itemize}
\item {} 
\sphinxAtStartPar
análisis de representaciones internas

\item {} 
\sphinxAtStartPar
regularización informacional

\item {} 
\sphinxAtStartPar
aprendizaje contrastivo

\end{itemize}


\bigskip\hrule\bigskip



\section{Desafíos Computacionales}
\label{\detokenize{03_multivariate_models/conditional_mi:desafios-computacionales}}\begin{itemize}
\item {} 
\sphinxAtStartPar
estimación en alta dimensión

\item {} 
\sphinxAtStartPar
sesgo en muestras pequeñas

\item {} 
\sphinxAtStartPar
necesidad de estimadores robustos

\item {} 
\sphinxAtStartPar
complejidad combinatoria

\end{itemize}

\sphinxAtStartPar
Métodos comunes:
\begin{itemize}
\item {} 
\sphinxAtStartPar
k\sphinxhyphen{}NN estimators

\item {} 
\sphinxAtStartPar
métodos kernel

\item {} 
\sphinxAtStartPar
modelos generativos

\item {} 
\sphinxAtStartPar
normalizing flows

\end{itemize}


\bigskip\hrule\bigskip



\section{Extensiones Avanzadas}
\label{\detokenize{03_multivariate_models/conditional_mi:extensiones-avanzadas}}\begin{itemize}
\item {} 
\sphinxAtStartPar
información mutua parcial

\item {} 
\sphinxAtStartPar
información multivariada total

\item {} 
\sphinxAtStartPar
sinergia e información redundante

\item {} 
\sphinxAtStartPar
descomposición PID (Partial Information Decomposition)

\end{itemize}


\bigskip\hrule\bigskip



\section{Conclusión}
\label{\detokenize{03_multivariate_models/conditional_mi:conclusion}}
\sphinxAtStartPar
La información mutua condicional multivariada permite analizar
dependencias directas eliminando la influencia de variables de control.

\sphinxAtStartPar
Es una herramienta esencial para:
\begin{itemize}
\item {} 
\sphinxAtStartPar
inferencia causal

\item {} 
\sphinxAtStartPar
modelamiento probabilístico

\item {} 
\sphinxAtStartPar
análisis estructural

\item {} 
\sphinxAtStartPar
aprendizaje automático moderno

\end{itemize}

\sphinxAtStartPar
En el siguiente capítulo se abordará el análisis de sistemas
de alta dimensión mediante herramientas computacionales
y notebooks experimentales.

\sphinxAtStartPar
{[}\hyperlink{cite.05_appendices/references:id9}{Pearl, 2009}{]}, {[}\hyperlink{cite.05_appendices/references:id5}{Murphy, 2012}{]}.

\sphinxstepscope


\chapter{Análisis de Información Mutua en Alta Dimensión}
\label{\detokenize{03_multivariate_models/high_dimensional:analisis-de-informacion-mutua-en-alta-dimension}}\label{\detokenize{03_multivariate_models/high_dimensional::doc}}
\sphinxAtStartPar
Este cuaderno presenta técnicas computacionales para analizar
dependencias informacionales en sistemas multivariados de alta dimensión.

\sphinxAtStartPar
Objetivos:
\begin{itemize}
\item {} 
\sphinxAtStartPar
Simular sistemas con dependencias estructuradas

\item {} 
\sphinxAtStartPar
Estimar información mutua entre múltiples variables

\item {} 
\sphinxAtStartPar
Construir matrices de información mutua

\item {} 
\sphinxAtStartPar
Analizar información mutua condicional

\item {} 
\sphinxAtStartPar
Evaluar desafíos computacionales en alta dimensión

\end{itemize}

\begin{sphinxuseclass}{cell}\begin{sphinxVerbatimInput}

\begin{sphinxuseclass}{cell_input}
\begin{sphinxVerbatim}[commandchars=\\\{\}]
\PYG{k+kn}{import}\PYG{+w}{ }\PYG{n+nn}{numpy}\PYG{+w}{ }\PYG{k}{as}\PYG{+w}{ }\PYG{n+nn}{np}
\PYG{k+kn}{import}\PYG{+w}{ }\PYG{n+nn}{pandas}\PYG{+w}{ }\PYG{k}{as}\PYG{+w}{ }\PYG{n+nn}{pd}
\PYG{k+kn}{import}\PYG{+w}{ }\PYG{n+nn}{matplotlib}\PYG{n+nn}{.}\PYG{n+nn}{pyplot}\PYG{+w}{ }\PYG{k}{as}\PYG{+w}{ }\PYG{n+nn}{plt}

\PYG{k+kn}{from}\PYG{+w}{ }\PYG{n+nn}{sklearn}\PYG{n+nn}{.}\PYG{n+nn}{feature\PYGZus{}selection}\PYG{+w}{ }\PYG{k+kn}{import} \PYG{n}{mutual\PYGZus{}info\PYGZus{}regression}
\PYG{k+kn}{from}\PYG{+w}{ }\PYG{n+nn}{sklearn}\PYG{n+nn}{.}\PYG{n+nn}{preprocessing}\PYG{+w}{ }\PYG{k+kn}{import} \PYG{n}{StandardScaler}
\end{sphinxVerbatim}

\end{sphinxuseclass}\end{sphinxVerbatimInput}

\end{sphinxuseclass}

\section{Generación de Datos Sintéticos Multivariados}
\label{\detokenize{03_multivariate_models/high_dimensional:generacion-de-datos-sinteticos-multivariados}}
\sphinxAtStartPar
Construiremos un sistema con:
\begin{itemize}
\item {} 
\sphinxAtStartPar
dependencias directas

\item {} 
\sphinxAtStartPar
relaciones no lineales

\item {} 
\sphinxAtStartPar
correlaciones indirectas

\end{itemize}

\begin{sphinxuseclass}{cell}\begin{sphinxVerbatimInput}

\begin{sphinxuseclass}{cell_input}
\begin{sphinxVerbatim}[commandchars=\\\{\}]
\PYG{n}{np}\PYG{o}{.}\PYG{n}{random}\PYG{o}{.}\PYG{n}{seed}\PYG{p}{(}\PYG{l+m+mi}{42}\PYG{p}{)}

\PYG{n}{n\PYGZus{}samples} \PYG{o}{=} \PYG{l+m+mi}{1500}

\PYG{n}{X1} \PYG{o}{=} \PYG{n}{np}\PYG{o}{.}\PYG{n}{random}\PYG{o}{.}\PYG{n}{normal}\PYG{p}{(}\PYG{n}{size}\PYG{o}{=}\PYG{n}{n\PYGZus{}samples}\PYG{p}{)}
\PYG{n}{X2} \PYG{o}{=} \PYG{l+m+mf}{0.8} \PYG{o}{*} \PYG{n}{X1} \PYG{o}{+} \PYG{n}{np}\PYG{o}{.}\PYG{n}{random}\PYG{o}{.}\PYG{n}{normal}\PYG{p}{(}\PYG{n}{scale}\PYG{o}{=}\PYG{l+m+mf}{0.4}\PYG{p}{,} \PYG{n}{size}\PYG{o}{=}\PYG{n}{n\PYGZus{}samples}\PYG{p}{)}
\PYG{n}{X3} \PYG{o}{=} \PYG{n}{np}\PYG{o}{.}\PYG{n}{sin}\PYG{p}{(}\PYG{n}{X1}\PYG{p}{)} \PYG{o}{+} \PYG{n}{np}\PYG{o}{.}\PYG{n}{random}\PYG{o}{.}\PYG{n}{normal}\PYG{p}{(}\PYG{n}{scale}\PYG{o}{=}\PYG{l+m+mf}{0.2}\PYG{p}{,} \PYG{n}{size}\PYG{o}{=}\PYG{n}{n\PYGZus{}samples}\PYG{p}{)}
\PYG{n}{X4} \PYG{o}{=} \PYG{n}{X2}\PYG{o}{*}\PYG{o}{*}\PYG{l+m+mi}{2} \PYG{o}{+} \PYG{n}{np}\PYG{o}{.}\PYG{n}{random}\PYG{o}{.}\PYG{n}{normal}\PYG{p}{(}\PYG{n}{scale}\PYG{o}{=}\PYG{l+m+mf}{0.3}\PYG{p}{,} \PYG{n}{size}\PYG{o}{=}\PYG{n}{n\PYGZus{}samples}\PYG{p}{)}
\PYG{n}{X5} \PYG{o}{=} \PYG{n}{np}\PYG{o}{.}\PYG{n}{random}\PYG{o}{.}\PYG{n}{normal}\PYG{p}{(}\PYG{n}{size}\PYG{o}{=}\PYG{n}{n\PYGZus{}samples}\PYG{p}{)}

\PYG{n}{data} \PYG{o}{=} \PYG{n}{np}\PYG{o}{.}\PYG{n}{vstack}\PYG{p}{(}\PYG{p}{[}\PYG{n}{X1}\PYG{p}{,} \PYG{n}{X2}\PYG{p}{,} \PYG{n}{X3}\PYG{p}{,} \PYG{n}{X4}\PYG{p}{,} \PYG{n}{X5}\PYG{p}{]}\PYG{p}{)}\PYG{o}{.}\PYG{n}{T}

\PYG{n}{df} \PYG{o}{=} \PYG{n}{pd}\PYG{o}{.}\PYG{n}{DataFrame}\PYG{p}{(}\PYG{n}{data}\PYG{p}{,} \PYG{n}{columns}\PYG{o}{=}\PYG{p}{[}\PYG{l+s+s2}{\PYGZdq{}}\PYG{l+s+s2}{X1}\PYG{l+s+s2}{\PYGZdq{}}\PYG{p}{,}\PYG{l+s+s2}{\PYGZdq{}}\PYG{l+s+s2}{X2}\PYG{l+s+s2}{\PYGZdq{}}\PYG{p}{,}\PYG{l+s+s2}{\PYGZdq{}}\PYG{l+s+s2}{X3}\PYG{l+s+s2}{\PYGZdq{}}\PYG{p}{,}\PYG{l+s+s2}{\PYGZdq{}}\PYG{l+s+s2}{X4}\PYG{l+s+s2}{\PYGZdq{}}\PYG{p}{,}\PYG{l+s+s2}{\PYGZdq{}}\PYG{l+s+s2}{X5}\PYG{l+s+s2}{\PYGZdq{}}\PYG{p}{]}\PYG{p}{)}

\PYG{n}{df}\PYG{o}{.}\PYG{n}{head}\PYG{p}{(}\PYG{p}{)}
\end{sphinxVerbatim}

\end{sphinxuseclass}\end{sphinxVerbatimInput}

\end{sphinxuseclass}

\section{Normalización}
\label{\detokenize{03_multivariate_models/high_dimensional:normalizacion}}
\sphinxAtStartPar
La normalización estabiliza estimaciones informacionales
en espacios de alta dimensión.

\begin{sphinxuseclass}{cell}\begin{sphinxVerbatimInput}

\begin{sphinxuseclass}{cell_input}
\begin{sphinxVerbatim}[commandchars=\\\{\}]
\PYG{n}{scaler} \PYG{o}{=} \PYG{n}{StandardScaler}\PYG{p}{(}\PYG{p}{)}
\PYG{n}{data\PYGZus{}scaled} \PYG{o}{=} \PYG{n}{scaler}\PYG{o}{.}\PYG{n}{fit\PYGZus{}transform}\PYG{p}{(}\PYG{n}{df}\PYG{o}{.}\PYG{n}{values}\PYG{p}{)}
\end{sphinxVerbatim}

\end{sphinxuseclass}\end{sphinxVerbatimInput}

\end{sphinxuseclass}

\section{Estimación de Matriz de Información Mutua}
\label{\detokenize{03_multivariate_models/high_dimensional:estimacion-de-matriz-de-informacion-mutua}}
\begin{sphinxuseclass}{cell}\begin{sphinxVerbatimInput}

\begin{sphinxuseclass}{cell_input}
\begin{sphinxVerbatim}[commandchars=\\\{\}]
\PYG{n}{n\PYGZus{}features} \PYG{o}{=} \PYG{n}{data\PYGZus{}scaled}\PYG{o}{.}\PYG{n}{shape}\PYG{p}{[}\PYG{l+m+mi}{1}\PYG{p}{]}
\PYG{n}{mi\PYGZus{}matrix} \PYG{o}{=} \PYG{n}{np}\PYG{o}{.}\PYG{n}{zeros}\PYG{p}{(}\PYG{p}{(}\PYG{n}{n\PYGZus{}features}\PYG{p}{,} \PYG{n}{n\PYGZus{}features}\PYG{p}{)}\PYG{p}{)}

\PYG{k}{for} \PYG{n}{i} \PYG{o+ow}{in} \PYG{n+nb}{range}\PYG{p}{(}\PYG{n}{n\PYGZus{}features}\PYG{p}{)}\PYG{p}{:}
    \PYG{k}{for} \PYG{n}{j} \PYG{o+ow}{in} \PYG{n+nb}{range}\PYG{p}{(}\PYG{n}{n\PYGZus{}features}\PYG{p}{)}\PYG{p}{:}
        \PYG{k}{if} \PYG{n}{i} \PYG{o}{!=} \PYG{n}{j}\PYG{p}{:}
            \PYG{n}{mi\PYGZus{}matrix}\PYG{p}{[}\PYG{n}{i}\PYG{p}{,} \PYG{n}{j}\PYG{p}{]} \PYG{o}{=} \PYG{n}{mutual\PYGZus{}info\PYGZus{}regression}\PYG{p}{(}
                \PYG{n}{data\PYGZus{}scaled}\PYG{p}{[}\PYG{p}{:}\PYG{p}{,} \PYG{p}{[}\PYG{n}{i}\PYG{p}{]}\PYG{p}{]}\PYG{p}{,}
                \PYG{n}{data\PYGZus{}scaled}\PYG{p}{[}\PYG{p}{:}\PYG{p}{,} \PYG{n}{j}\PYG{p}{]}\PYG{p}{,}
                \PYG{n}{random\PYGZus{}state}\PYG{o}{=}\PYG{l+m+mi}{42}
            \PYG{p}{)}\PYG{p}{[}\PYG{l+m+mi}{0}\PYG{p}{]}

\PYG{n}{mi\PYGZus{}df} \PYG{o}{=} \PYG{n}{pd}\PYG{o}{.}\PYG{n}{DataFrame}\PYG{p}{(}\PYG{n}{mi\PYGZus{}matrix}\PYG{p}{,}
                     \PYG{n}{columns}\PYG{o}{=}\PYG{n}{df}\PYG{o}{.}\PYG{n}{columns}\PYG{p}{,}
                     \PYG{n}{index}\PYG{o}{=}\PYG{n}{df}\PYG{o}{.}\PYG{n}{columns}\PYG{p}{)}

\PYG{n}{mi\PYGZus{}df}
\end{sphinxVerbatim}

\end{sphinxuseclass}\end{sphinxVerbatimInput}

\end{sphinxuseclass}

\section{Visualización de la Matriz MI}
\label{\detokenize{03_multivariate_models/high_dimensional:visualizacion-de-la-matriz-mi}}
\begin{sphinxuseclass}{cell}\begin{sphinxVerbatimInput}

\begin{sphinxuseclass}{cell_input}
\begin{sphinxVerbatim}[commandchars=\\\{\}]
\PYG{n}{plt}\PYG{o}{.}\PYG{n}{figure}\PYG{p}{(}\PYG{p}{)}

\PYG{n}{plt}\PYG{o}{.}\PYG{n}{imshow}\PYG{p}{(}\PYG{n}{mi\PYGZus{}matrix}\PYG{p}{)}
\PYG{n}{plt}\PYG{o}{.}\PYG{n}{colorbar}\PYG{p}{(}\PYG{p}{)}

\PYG{n}{plt}\PYG{o}{.}\PYG{n}{xticks}\PYG{p}{(}\PYG{n+nb}{range}\PYG{p}{(}\PYG{n}{n\PYGZus{}features}\PYG{p}{)}\PYG{p}{,} \PYG{n}{df}\PYG{o}{.}\PYG{n}{columns}\PYG{p}{)}
\PYG{n}{plt}\PYG{o}{.}\PYG{n}{yticks}\PYG{p}{(}\PYG{n+nb}{range}\PYG{p}{(}\PYG{n}{n\PYGZus{}features}\PYG{p}{)}\PYG{p}{,} \PYG{n}{df}\PYG{o}{.}\PYG{n}{columns}\PYG{p}{)}

\PYG{n}{plt}\PYG{o}{.}\PYG{n}{title}\PYG{p}{(}\PYG{l+s+s2}{\PYGZdq{}}\PYG{l+s+s2}{Matriz de Información Mutua}\PYG{l+s+s2}{\PYGZdq{}}\PYG{p}{)}
\PYG{n}{plt}\PYG{o}{.}\PYG{n}{show}\PYG{p}{(}\PYG{p}{)}
\end{sphinxVerbatim}

\end{sphinxuseclass}\end{sphinxVerbatimInput}

\end{sphinxuseclass}

\section{Información Mutua Condicional Aproximada}
\label{\detokenize{03_multivariate_models/high_dimensional:informacion-mutua-condicional-aproximada}}
\sphinxAtStartPar
Aproximaremos:

\sphinxAtStartPar
I(Xi ; Xj | Xk)

\sphinxAtStartPar
mediante regresión residual.

\begin{sphinxuseclass}{cell}\begin{sphinxVerbatimInput}

\begin{sphinxuseclass}{cell_input}
\begin{sphinxVerbatim}[commandchars=\\\{\}]
\PYG{k+kn}{from}\PYG{+w}{ }\PYG{n+nn}{sklearn}\PYG{n+nn}{.}\PYG{n+nn}{linear\PYGZus{}model}\PYG{+w}{ }\PYG{k+kn}{import} \PYG{n}{LinearRegression}

\PYG{k}{def}\PYG{+w}{ }\PYG{n+nf}{conditional\PYGZus{}mi\PYGZus{}residual}\PYG{p}{(}\PYG{n}{X}\PYG{p}{,} \PYG{n}{Y}\PYG{p}{,} \PYG{n}{Z}\PYG{p}{)}\PYG{p}{:}
    \PYG{n}{model\PYGZus{}x} \PYG{o}{=} \PYG{n}{LinearRegression}\PYG{p}{(}\PYG{p}{)}\PYG{o}{.}\PYG{n}{fit}\PYG{p}{(}\PYG{n}{Z}\PYG{p}{,} \PYG{n}{X}\PYG{p}{)}
    \PYG{n}{model\PYGZus{}y} \PYG{o}{=} \PYG{n}{LinearRegression}\PYG{p}{(}\PYG{p}{)}\PYG{o}{.}\PYG{n}{fit}\PYG{p}{(}\PYG{n}{Z}\PYG{p}{,} \PYG{n}{Y}\PYG{p}{)}

    \PYG{n}{rx} \PYG{o}{=} \PYG{n}{X} \PYG{o}{\PYGZhy{}} \PYG{n}{model\PYGZus{}x}\PYG{o}{.}\PYG{n}{predict}\PYG{p}{(}\PYG{n}{Z}\PYG{p}{)}
    \PYG{n}{ry} \PYG{o}{=} \PYG{n}{Y} \PYG{o}{\PYGZhy{}} \PYG{n}{model\PYGZus{}y}\PYG{o}{.}\PYG{n}{predict}\PYG{p}{(}\PYG{n}{Z}\PYG{p}{)}

    \PYG{k}{return} \PYG{n}{mutual\PYGZus{}info\PYGZus{}regression}\PYG{p}{(}
        \PYG{n}{rx}\PYG{o}{.}\PYG{n}{reshape}\PYG{p}{(}\PYG{o}{\PYGZhy{}}\PYG{l+m+mi}{1}\PYG{p}{,}\PYG{l+m+mi}{1}\PYG{p}{)}\PYG{p}{,}
        \PYG{n}{ry}\PYG{p}{,}
        \PYG{n}{random\PYGZus{}state}\PYG{o}{=}\PYG{l+m+mi}{42}
    \PYG{p}{)}\PYG{p}{[}\PYG{l+m+mi}{0}\PYG{p}{]}
\end{sphinxVerbatim}

\end{sphinxuseclass}\end{sphinxVerbatimInput}

\end{sphinxuseclass}
\begin{sphinxuseclass}{cell}\begin{sphinxVerbatimInput}

\begin{sphinxuseclass}{cell_input}
\begin{sphinxVerbatim}[commandchars=\\\{\}]
\PYG{n}{cmiv} \PYG{o}{=} \PYG{n}{conditional\PYGZus{}mi\PYGZus{}residual}\PYG{p}{(}
    \PYG{n}{data\PYGZus{}scaled}\PYG{p}{[}\PYG{p}{:}\PYG{p}{,}\PYG{l+m+mi}{0}\PYG{p}{]}\PYG{p}{,}
    \PYG{n}{data\PYGZus{}scaled}\PYG{p}{[}\PYG{p}{:}\PYG{p}{,}\PYG{l+m+mi}{3}\PYG{p}{]}\PYG{p}{,}
    \PYG{n}{data\PYGZus{}scaled}\PYG{p}{[}\PYG{p}{:}\PYG{p}{,}\PYG{l+m+mi}{1}\PYG{p}{]}\PYG{o}{.}\PYG{n}{reshape}\PYG{p}{(}\PYG{o}{\PYGZhy{}}\PYG{l+m+mi}{1}\PYG{p}{,}\PYG{l+m+mi}{1}\PYG{p}{)}
\PYG{p}{)}

\PYG{n+nb}{print}\PYG{p}{(}\PYG{l+s+s2}{\PYGZdq{}}\PYG{l+s+s2}{I(X1 ; X4 | X2) ≈}\PYG{l+s+s2}{\PYGZdq{}}\PYG{p}{,} \PYG{n}{cmiv}\PYG{p}{)}
\end{sphinxVerbatim}

\end{sphinxuseclass}\end{sphinxVerbatimInput}

\end{sphinxuseclass}

\section{Crecimiento de Dimensionalidad}
\label{\detokenize{03_multivariate_models/high_dimensional:crecimiento-de-dimensionalidad}}
\sphinxAtStartPar
Analizamos el comportamiento de la estimación
cuando aumenta el número de variables.

\begin{sphinxuseclass}{cell}\begin{sphinxVerbatimInput}

\begin{sphinxuseclass}{cell_input}
\begin{sphinxVerbatim}[commandchars=\\\{\}]
\PYG{n}{dims} \PYG{o}{=} \PYG{p}{[}\PYG{p}{]}
\PYG{n}{times} \PYG{o}{=} \PYG{p}{[}\PYG{p}{]}

\PYG{k}{for} \PYG{n}{d} \PYG{o+ow}{in} \PYG{p}{[}\PYG{l+m+mi}{5}\PYG{p}{,}\PYG{l+m+mi}{10}\PYG{p}{,}\PYG{l+m+mi}{20}\PYG{p}{,}\PYG{l+m+mi}{30}\PYG{p}{]}\PYG{p}{:}
    \PYG{n}{X} \PYG{o}{=} \PYG{n}{np}\PYG{o}{.}\PYG{n}{random}\PYG{o}{.}\PYG{n}{normal}\PYG{p}{(}\PYG{n}{size}\PYG{o}{=}\PYG{p}{(}\PYG{l+m+mi}{800}\PYG{p}{,}\PYG{n}{d}\PYG{p}{)}\PYG{p}{)}

    \PYG{k+kn}{import}\PYG{+w}{ }\PYG{n+nn}{time}
    \PYG{n}{t0} \PYG{o}{=} \PYG{n}{time}\PYG{o}{.}\PYG{n}{time}\PYG{p}{(}\PYG{p}{)}

    \PYG{k}{for} \PYG{n}{i} \PYG{o+ow}{in} \PYG{n+nb}{range}\PYG{p}{(}\PYG{n}{d}\PYG{p}{)}\PYG{p}{:}
        \PYG{k}{for} \PYG{n}{j} \PYG{o+ow}{in} \PYG{n+nb}{range}\PYG{p}{(}\PYG{n}{d}\PYG{p}{)}\PYG{p}{:}
            \PYG{k}{if} \PYG{n}{i} \PYG{o}{!=} \PYG{n}{j}\PYG{p}{:}
                \PYG{n}{mutual\PYGZus{}info\PYGZus{}regression}\PYG{p}{(}\PYG{n}{X}\PYG{p}{[}\PYG{p}{:}\PYG{p}{,}\PYG{p}{[}\PYG{n}{i}\PYG{p}{]}\PYG{p}{]}\PYG{p}{,} \PYG{n}{X}\PYG{p}{[}\PYG{p}{:}\PYG{p}{,}\PYG{n}{j}\PYG{p}{]}\PYG{p}{)}

    \PYG{n}{t1} \PYG{o}{=} \PYG{n}{time}\PYG{o}{.}\PYG{n}{time}\PYG{p}{(}\PYG{p}{)}

    \PYG{n}{dims}\PYG{o}{.}\PYG{n}{append}\PYG{p}{(}\PYG{n}{d}\PYG{p}{)}
    \PYG{n}{times}\PYG{o}{.}\PYG{n}{append}\PYG{p}{(}\PYG{n}{t1} \PYG{o}{\PYGZhy{}} \PYG{n}{t0}\PYG{p}{)}
\end{sphinxVerbatim}

\end{sphinxuseclass}\end{sphinxVerbatimInput}

\end{sphinxuseclass}
\begin{sphinxuseclass}{cell}\begin{sphinxVerbatimInput}

\begin{sphinxuseclass}{cell_input}
\begin{sphinxVerbatim}[commandchars=\\\{\}]
\PYG{n}{plt}\PYG{o}{.}\PYG{n}{figure}\PYG{p}{(}\PYG{p}{)}
\PYG{n}{plt}\PYG{o}{.}\PYG{n}{plot}\PYG{p}{(}\PYG{n}{dims}\PYG{p}{,} \PYG{n}{times}\PYG{p}{,} \PYG{n}{marker}\PYG{o}{=}\PYG{l+s+s2}{\PYGZdq{}}\PYG{l+s+s2}{o}\PYG{l+s+s2}{\PYGZdq{}}\PYG{p}{)}

\PYG{n}{plt}\PYG{o}{.}\PYG{n}{xlabel}\PYG{p}{(}\PYG{l+s+s2}{\PYGZdq{}}\PYG{l+s+s2}{Dimensión}\PYG{l+s+s2}{\PYGZdq{}}\PYG{p}{)}
\PYG{n}{plt}\PYG{o}{.}\PYG{n}{ylabel}\PYG{p}{(}\PYG{l+s+s2}{\PYGZdq{}}\PYG{l+s+s2}{Tiempo (s)}\PYG{l+s+s2}{\PYGZdq{}}\PYG{p}{)}
\PYG{n}{plt}\PYG{o}{.}\PYG{n}{title}\PYG{p}{(}\PYG{l+s+s2}{\PYGZdq{}}\PYG{l+s+s2}{Costo Computacional en Alta Dimensión}\PYG{l+s+s2}{\PYGZdq{}}\PYG{p}{)}

\PYG{n}{plt}\PYG{o}{.}\PYG{n}{show}\PYG{p}{(}\PYG{p}{)}
\end{sphinxVerbatim}

\end{sphinxuseclass}\end{sphinxVerbatimInput}

\end{sphinxuseclass}

\section{Discusión}
\label{\detokenize{03_multivariate_models/high_dimensional:discusion}}
\sphinxAtStartPar
Observaciones principales:
\begin{itemize}
\item {} 
\sphinxAtStartPar
La información mutua detecta dependencias no lineales.

\item {} 
\sphinxAtStartPar
La dimensionalidad incrementa rápidamente el costo computacional.

\item {} 
\sphinxAtStartPar
La información mutua condicional permite eliminar dependencias indirectas.

\item {} 
\sphinxAtStartPar
La estimación robusta en alta dimensión requiere métodos avanzados.

\end{itemize}

\sphinxAtStartPar
Este análisis establece las bases para aplicaciones reales
en ciencia de datos y aprendizaje profundo.

\sphinxAtStartPar
{[}\hyperlink{cite.05_appendices/references:id7}{Kraskov \sphinxstyleemphasis{et al.}, 2004}{]}.

\sphinxstepscope


\part{Aplicaciones en Ciencia de Datos}

\sphinxstepscope


\chapter{Clustering basado en Información Mutua}
\label{\detokenize{04_data_science_applications/clustering_mi:clustering-basado-en-informacion-mutua}}\label{\detokenize{04_data_science_applications/clustering_mi::doc}}

\section{Introducción}
\label{\detokenize{04_data_science_applications/clustering_mi:introduccion}}
\sphinxAtStartPar
La información mutua constituye una herramienta fundamental para medir la dependencia estadística entre variables aleatorias.\\
En ciencia de datos, se utiliza ampliamente para:
\begin{itemize}
\item {} 
\sphinxAtStartPar
Evaluar calidad de agrupamientos

\item {} 
\sphinxAtStartPar
Diseñar funciones objetivo

\item {} 
\sphinxAtStartPar
Comparar particiones

\item {} 
\sphinxAtStartPar
Detectar estructura latente en datos

\end{itemize}

\sphinxAtStartPar
A diferencia de métricas basadas en distancia euclidiana, la información mutua captura relaciones no lineales y dependencias complejas.


\bigskip\hrule\bigskip



\section{Definición de Información Mutua para Clustering}
\label{\detokenize{04_data_science_applications/clustering_mi:definicion-de-informacion-mutua-para-clustering}}
\sphinxAtStartPar
Sean:
\begin{itemize}
\item {} 
\sphinxAtStartPar
\(C\) : etiquetas reales

\item {} 
\sphinxAtStartPar
\(\hat{C}\) : etiquetas obtenidas por el algoritmo de clustering

\end{itemize}

\sphinxAtStartPar
La información mutua se define como:
\begin{equation*}
\begin{split}
I(C;\hat{C}) = \sum_{c,\hat{c}} P(c,\hat{c})
\log \frac{P(c,\hat{c})}{P(c)P(\hat{c})}
\end{split}
\end{equation*}
\sphinxAtStartPar
Interpretación:
\begin{itemize}
\item {} 
\sphinxAtStartPar
\(I=0\) → independencia total

\item {} 
\sphinxAtStartPar
\(I\) grande → fuerte concordancia estructural

\end{itemize}


\bigskip\hrule\bigskip



\section{Normalized Mutual Information (NMI)}
\label{\detokenize{04_data_science_applications/clustering_mi:normalized-mutual-information-nmi}}
\sphinxAtStartPar
Para comparar clustering con distinto número de grupos se utiliza:
\begin{equation*}
\begin{split}
NMI = \frac{I(C;\hat{C})}{\sqrt{H(C)H(\hat{C})}}
\end{split}
\end{equation*}
\sphinxAtStartPar
Propiedades:
\begin{itemize}
\item {} 
\sphinxAtStartPar
\(0 \leq NMI \leq 1\)

\item {} 
\sphinxAtStartPar
Invariante ante permutaciones de etiquetas

\item {} 
\sphinxAtStartPar
Robusta frente a clusters desbalanceados

\end{itemize}


\bigskip\hrule\bigskip



\section{Ajuste por Azar: Adjusted Mutual Information (AMI)}
\label{\detokenize{04_data_science_applications/clustering_mi:ajuste-por-azar-adjusted-mutual-information-ami}}
\sphinxAtStartPar
Corrige el sesgo introducido por coincidencias aleatorias:
\begin{equation*}
\begin{split}
AMI = \frac{I - E[I]}{\max(H(C),H(\hat{C})) - E[I]}
\end{split}
\end{equation*}
\sphinxAtStartPar
Ventajas:
\begin{itemize}
\item {} 
\sphinxAtStartPar
Penaliza coincidencias fortuitas

\item {} 
\sphinxAtStartPar
Mejora comparabilidad entre modelos

\end{itemize}


\bigskip\hrule\bigskip



\section{Ejemplo Conceptual}
\label{\detokenize{04_data_science_applications/clustering_mi:ejemplo-conceptual}}
\sphinxAtStartPar
Supongamos:
\begin{itemize}
\item {} 
\sphinxAtStartPar
Clases reales: 3

\item {} 
\sphinxAtStartPar
Clusters obtenidos: 3

\end{itemize}

\sphinxAtStartPar
Matriz de contingencia:


\begin{savenotes}\sphinxattablestart
\centering
\begin{tabulary}{\linewidth}[t]{|T|T|T|T|}
\hline
\sphinxstyletheadfamily 
\sphinxAtStartPar
Real / Pred
&\sphinxstyletheadfamily 
\sphinxAtStartPar
A
&\sphinxstyletheadfamily 
\sphinxAtStartPar
B
&\sphinxstyletheadfamily 
\sphinxAtStartPar
C
\\
\hline
\sphinxAtStartPar
1
&
\sphinxAtStartPar
30
&
\sphinxAtStartPar
2
&
\sphinxAtStartPar
1
\\
\hline
\sphinxAtStartPar
2
&
\sphinxAtStartPar
3
&
\sphinxAtStartPar
25
&
\sphinxAtStartPar
2
\\
\hline
\sphinxAtStartPar
3
&
\sphinxAtStartPar
1
&
\sphinxAtStartPar
4
&
\sphinxAtStartPar
28
\\
\hline
\end{tabulary}
\par
\sphinxattableend\end{savenotes}

\sphinxAtStartPar
A partir de esta tabla:
\begin{enumerate}
\sphinxsetlistlabels{\arabic}{enumi}{enumii}{}{.}%
\item {} 
\sphinxAtStartPar
Se calcula \(P(c,\hat{c})\)

\item {} 
\sphinxAtStartPar
Se obtienen marginales

\item {} 
\sphinxAtStartPar
Se evalúa la suma de información mutua

\end{enumerate}

\sphinxAtStartPar
Resultado esperado:
\begin{itemize}
\item {} 
\sphinxAtStartPar
NMI alto

\item {} 
\sphinxAtStartPar
Buen alineamiento estructural

\end{itemize}


\bigskip\hrule\bigskip



\section{Información Mutua como Función Objetivo}
\label{\detokenize{04_data_science_applications/clustering_mi:informacion-mutua-como-funcion-objetivo}}
\sphinxAtStartPar
Algunos algoritmos maximizan directamente:
\begin{equation*}
\begin{split}
\max I(X;Z)
\end{split}
\end{equation*}
\sphinxAtStartPar
donde:
\begin{itemize}
\item {} 
\sphinxAtStartPar
\(X\) datos originales

\item {} 
\sphinxAtStartPar
\(Z\) representación latente

\end{itemize}

\sphinxAtStartPar
Ejemplos:
\begin{itemize}
\item {} 
\sphinxAtStartPar
Information Bottleneck Clustering

\item {} 
\sphinxAtStartPar
Deep InfoMax Clustering

\item {} 
\sphinxAtStartPar
Spectral Clustering basado en MI

\end{itemize}


\bigskip\hrule\bigskip



\section{Ventajas frente a Distancias Clásicas}
\label{\detokenize{04_data_science_applications/clustering_mi:ventajas-frente-a-distancias-clasicas}}

\begin{savenotes}\sphinxattablestart
\centering
\begin{tabulary}{\linewidth}[t]{|T|T|T|T|}
\hline
\sphinxstyletheadfamily 
\sphinxAtStartPar
Métrica
&\sphinxstyletheadfamily 
\sphinxAtStartPar
Dependencias no lineales
&\sphinxstyletheadfamily 
\sphinxAtStartPar
Invarianza
&\sphinxstyletheadfamily 
\sphinxAtStartPar
Robustez
\\
\hline
\sphinxAtStartPar
Euclidiana
&
\sphinxAtStartPar
No
&
\sphinxAtStartPar
Baja
&
\sphinxAtStartPar
Media
\\
\hline
\sphinxAtStartPar
Coseno
&
\sphinxAtStartPar
Parcial
&
\sphinxAtStartPar
Media
&
\sphinxAtStartPar
Media
\\
\hline
\sphinxAtStartPar
Información Mutua
&
\sphinxAtStartPar
Sí
&
\sphinxAtStartPar
Alta
&
\sphinxAtStartPar
Alta
\\
\hline
\end{tabulary}
\par
\sphinxattableend\end{savenotes}


\bigskip\hrule\bigskip



\section{Limitaciones}
\label{\detokenize{04_data_science_applications/clustering_mi:limitaciones}}\begin{itemize}
\item {} 
\sphinxAtStartPar
Estimación de probabilidades en alta dimensión

\item {} 
\sphinxAtStartPar
Sensibilidad a discretización

\item {} 
\sphinxAtStartPar
Costo computacional

\end{itemize}

\sphinxAtStartPar
Soluciones modernas:
\begin{itemize}
\item {} 
\sphinxAtStartPar
Estimadores kNN

\item {} 
\sphinxAtStartPar
KDE

\item {} 
\sphinxAtStartPar
Redes neuronales estimadoras de MI

\end{itemize}


\bigskip\hrule\bigskip



\section{Conclusiones}
\label{\detokenize{04_data_science_applications/clustering_mi:conclusiones}}
\sphinxAtStartPar
La información mutua permite:
\begin{itemize}
\item {} 
\sphinxAtStartPar
Evaluar clustering sin suposiciones geométricas

\item {} 
\sphinxAtStartPar
Capturar relaciones complejas

\item {} 
\sphinxAtStartPar
Diseñar algoritmos robustos

\end{itemize}

\sphinxAtStartPar
Es especialmente útil en:
\begin{itemize}
\item {} 
\sphinxAtStartPar
Datos no lineales

\item {} 
\sphinxAtStartPar
Variables categóricas

\item {} 
\sphinxAtStartPar
Representaciones profundas

\end{itemize}

\sphinxAtStartPar
{[}\hyperlink{cite.05_appendices/references:id4}{Bishop, 2006}{]}, {[}\hyperlink{cite.05_appendices/references:id5}{Murphy, 2012}{]}.

\sphinxstepscope


\chapter{Información Mutua en Deep Learning}
\label{\detokenize{04_data_science_applications/deep_learning_mi:informacion-mutua-en-deep-learning}}\label{\detokenize{04_data_science_applications/deep_learning_mi::doc}}

\section{Introducción}
\label{\detokenize{04_data_science_applications/deep_learning_mi:introduccion}}
\sphinxAtStartPar
La información mutua ha emergido como una herramienta central en el análisis y diseño de modelos de aprendizaje profundo.\\
Permite cuantificar:
\begin{itemize}
\item {} 
\sphinxAtStartPar
Cantidad de información preservada

\item {} 
\sphinxAtStartPar
Compresión de representaciones

\item {} 
\sphinxAtStartPar
Dependencia entre capas

\end{itemize}

\sphinxAtStartPar
Se aplica en:
\begin{itemize}
\item {} 
\sphinxAtStartPar
aprendizaje auto\sphinxhyphen{}supervisado

\item {} 
\sphinxAtStartPar
aprendizaje contrastivo

\item {} 
\sphinxAtStartPar
representación latente

\item {} 
\sphinxAtStartPar
regularización informacional

\end{itemize}


\bigskip\hrule\bigskip



\section{Marco Teórico: Information Bottleneck}
\label{\detokenize{04_data_science_applications/deep_learning_mi:marco-teorico-information-bottleneck}}
\sphinxAtStartPar
Sea:
\begin{itemize}
\item {} 
\sphinxAtStartPar
\(X\) entrada

\item {} 
\sphinxAtStartPar
\(Y\) salida

\item {} 
\sphinxAtStartPar
\(Z\) representación intermedia

\end{itemize}

\sphinxAtStartPar
Objetivo:
\begin{equation*}
\begin{split}
\min I(X;Z) - \beta I(Z;Y)
\end{split}
\end{equation*}
\sphinxAtStartPar
Interpretación:
\begin{itemize}
\item {} 
\sphinxAtStartPar
minimizar redundancia

\item {} 
\sphinxAtStartPar
maximizar información relevante

\end{itemize}


\bigskip\hrule\bigskip



\section{Deep InfoMax}
\label{\detokenize{04_data_science_applications/deep_learning_mi:deep-infomax}}
\sphinxAtStartPar
Busca maximizar:
\begin{equation*}
\begin{split}
I(X;Z)
\end{split}
\end{equation*}
\sphinxAtStartPar
entre:
\begin{itemize}
\item {} 
\sphinxAtStartPar
entrada global

\item {} 
\sphinxAtStartPar
representación latente

\end{itemize}

\sphinxAtStartPar
Aplicaciones:
\begin{itemize}
\item {} 
\sphinxAtStartPar
aprendizaje no supervisado

\item {} 
\sphinxAtStartPar
representación robusta

\item {} 
\sphinxAtStartPar
embeddings semánticos

\end{itemize}


\bigskip\hrule\bigskip



\section{Aprendizaje Contrastivo}
\label{\detokenize{04_data_science_applications/deep_learning_mi:aprendizaje-contrastivo}}
\sphinxAtStartPar
Métodos como:
\begin{itemize}
\item {} 
\sphinxAtStartPar
SimCLR

\item {} 
\sphinxAtStartPar
MoCo

\item {} 
\sphinxAtStartPar
BYOL

\item {} 
\sphinxAtStartPar
InfoNCE

\end{itemize}

\sphinxAtStartPar
optimizan aproximaciones de información mutua.

\sphinxAtStartPar
Función InfoNCE:
\begin{equation*}
\begin{split}
\mathcal{L} =
-\mathbb{E}
\left[
\log
\frac{\exp(f(x,x^+))}
{\sum_j \exp(f(x,x_j))}
\right]
\end{split}
\end{equation*}
\sphinxAtStartPar
Interpretación:
\begin{itemize}
\item {} 
\sphinxAtStartPar
maximiza similitud positiva

\item {} 
\sphinxAtStartPar
minimiza similitud negativa

\end{itemize}


\bigskip\hrule\bigskip



\section{Estimadores Neuronales de Información Mutua}
\label{\detokenize{04_data_science_applications/deep_learning_mi:estimadores-neuronales-de-informacion-mutua}}
\sphinxAtStartPar
Debido a la dificultad de estimar MI directamente:


\subsection{MINE (Mutual Information Neural Estimation)}
\label{\detokenize{04_data_science_applications/deep_learning_mi:mine-mutual-information-neural-estimation}}\begin{equation*}
\begin{split}
I(X;Y) \ge
E[T_\theta] -
\log E[e^{T_\theta}]
\end{split}
\end{equation*}
\sphinxAtStartPar
Ventajas:
\begin{itemize}
\item {} 
\sphinxAtStartPar
escalable

\item {} 
\sphinxAtStartPar
diferenciable

\item {} 
\sphinxAtStartPar
apto para entrenamiento end\sphinxhyphen{}to\sphinxhyphen{}end

\end{itemize}


\bigskip\hrule\bigskip



\section{Aplicaciones en Redes Profundas}
\label{\detokenize{04_data_science_applications/deep_learning_mi:aplicaciones-en-redes-profundas}}

\subsection{1. Autoencoders Variacionales}
\label{\detokenize{04_data_science_applications/deep_learning_mi:autoencoders-variacionales}}
\sphinxAtStartPar
Relación con divergencia KL:
\begin{equation*}
\begin{split}
KL(q(z|x) || p(z))
\end{split}
\end{equation*}
\sphinxAtStartPar
Controla:
\begin{itemize}
\item {} 
\sphinxAtStartPar
regularización

\item {} 
\sphinxAtStartPar
compresión latente

\end{itemize}


\bigskip\hrule\bigskip



\subsection{2. Representación Auto\sphinxhyphen{}supervisada}
\label{\detokenize{04_data_science_applications/deep_learning_mi:representacion-auto-supervisada}}
\sphinxAtStartPar
Maximizar MI entre:
\begin{itemize}
\item {} 
\sphinxAtStartPar
vistas aumentadas

\item {} 
\sphinxAtStartPar
proyecciones latentes

\end{itemize}

\sphinxAtStartPar
Resultados:
\begin{itemize}
\item {} 
\sphinxAtStartPar
robustez

\item {} 
\sphinxAtStartPar
generalización

\item {} 
\sphinxAtStartPar
invariancia

\end{itemize}


\bigskip\hrule\bigskip



\subsection{3. Redes Generativas}
\label{\detokenize{04_data_science_applications/deep_learning_mi:redes-generativas}}
\sphinxAtStartPar
GANs y modelos generativos usan MI para:
\begin{itemize}
\item {} 
\sphinxAtStartPar
disentanglement

\item {} 
\sphinxAtStartPar
control semántico

\item {} 
\sphinxAtStartPar
interpretabilidad

\end{itemize}

\sphinxAtStartPar
Ejemplo: InfoGAN
\begin{equation*}
\begin{split}
\max I(c;G(z,c))
\end{split}
\end{equation*}

\bigskip\hrule\bigskip



\section{Limitaciones Prácticas}
\label{\detokenize{04_data_science_applications/deep_learning_mi:limitaciones-practicas}}\begin{itemize}
\item {} 
\sphinxAtStartPar
estimación sesgada

\item {} 
\sphinxAtStartPar
explosión numérica

\item {} 
\sphinxAtStartPar
alta varianza

\end{itemize}

\sphinxAtStartPar
Soluciones:
\begin{itemize}
\item {} 
\sphinxAtStartPar
estimadores contrastivos

\item {} 
\sphinxAtStartPar
regularización espectral

\item {} 
\sphinxAtStartPar
clipping de gradientes

\end{itemize}


\bigskip\hrule\bigskip



\section{Tendencias Actuales en Investigación}
\label{\detokenize{04_data_science_applications/deep_learning_mi:tendencias-actuales-en-investigacion}}\begin{itemize}
\item {} 
\sphinxAtStartPar
Self\sphinxhyphen{}supervised learning

\item {} 
\sphinxAtStartPar
Multimodal representation learning

\item {} 
\sphinxAtStartPar
Graph neural networks informacionales

\item {} 
\sphinxAtStartPar
Causal representation learning

\end{itemize}


\bigskip\hrule\bigskip



\section{Conclusiones}
\label{\detokenize{04_data_science_applications/deep_learning_mi:conclusiones}}
\sphinxAtStartPar
La información mutua permite:
\begin{itemize}
\item {} 
\sphinxAtStartPar
interpretar redes profundas

\item {} 
\sphinxAtStartPar
diseñar funciones objetivo informacionales

\item {} 
\sphinxAtStartPar
mejorar representaciones latentes

\end{itemize}

\sphinxAtStartPar
Se ha convertido en un componente clave del deep learning moderno.

\sphinxAtStartPar
{[}\hyperlink{cite.05_appendices/references:id8}{Goodfellow \sphinxstyleemphasis{et al.}, 2016}{]}.

\sphinxstepscope


\part{Apéndices}

\sphinxstepscope


\chapter{Demostraciones Matemáticas Avanzadas}
\label{\detokenize{05_appendices/proofs:demostraciones-matematicas-avanzadas}}\label{\detokenize{05_appendices/proofs::doc}}

\section{Introducción}
\label{\detokenize{05_appendices/proofs:introduccion}}
\sphinxAtStartPar
Este apéndice presenta demostraciones formales de resultados fundamentales de la teoría de la información utilizados a lo largo del libro.\\
El enfoque combina rigor matemático con interpretaciones informacionales relevantes para ciencia de datos e ingeniería.


\bigskip\hrule\bigskip



\section{A.1 No Negatividad de la Información Mutua}
\label{\detokenize{05_appendices/proofs:a-1-no-negatividad-de-la-informacion-mutua}}

\subsection{Teorema}
\label{\detokenize{05_appendices/proofs:teorema}}
\sphinxAtStartPar
Para variables aleatorias discretas \(X\) e \(Y\):
\begin{equation*}
\begin{split}
I(X;Y) \ge 0
\end{split}
\end{equation*}
\sphinxAtStartPar
con igualdad si y sólo si \(X\) e \(Y\) son independientes.


\bigskip\hrule\bigskip



\subsection{Demostración}
\label{\detokenize{05_appendices/proofs:demostracion}}
\sphinxAtStartPar
Por definición:
\begin{equation*}
\begin{split}
I(X;Y)=
\sum_{x,y} P(x,y)
\log \frac{P(x,y)}{P(x)P(y)}
\end{split}
\end{equation*}
\sphinxAtStartPar
Esto corresponde exactamente a la divergencia KL:
\begin{equation*}
\begin{split}
I(X;Y)=
D_{KL}\big(P(X,Y)\parallel P(X)P(Y)\big)
\end{split}
\end{equation*}
\sphinxAtStartPar
Sabemos que:
\begin{equation*}
\begin{split}
D_{KL}(P||Q)\ge 0
\end{split}
\end{equation*}
\sphinxAtStartPar
por la desigualdad de Gibbs.

\sphinxAtStartPar
Por lo tanto:
\begin{equation*}
\begin{split}
I(X;Y)\ge0
\end{split}
\end{equation*}
\sphinxAtStartPar
La igualdad ocurre únicamente cuando:
\begin{equation*}
\begin{split}
P(x,y)=P(x)P(y)
\end{split}
\end{equation*}
\sphinxAtStartPar
lo que implica independencia estadística.


\bigskip\hrule\bigskip



\section{A.2 Regla de la Cadena para la Entropía}
\label{\detokenize{05_appendices/proofs:a-2-regla-de-la-cadena-para-la-entropia}}

\subsection{Teorema}
\label{\detokenize{05_appendices/proofs:id1}}\begin{equation*}
\begin{split}
H(X,Y)=H(X)+H(Y|X)
\end{split}
\end{equation*}

\bigskip\hrule\bigskip



\subsection{Demostración}
\label{\detokenize{05_appendices/proofs:id2}}
\sphinxAtStartPar
Definición de entropía conjunta:
\begin{equation*}
\begin{split}
H(X,Y)=-\sum_{x,y}P(x,y)\log P(x,y)
\end{split}
\end{equation*}
\sphinxAtStartPar
Usamos:
\begin{equation*}
\begin{split}
P(x,y)=P(x)P(y|x)
\end{split}
\end{equation*}
\sphinxAtStartPar
entonces:
\begin{equation*}
\begin{split}
\log P(x,y)=\log P(x)+\log P(y|x)
\end{split}
\end{equation*}
\sphinxAtStartPar
Sustituyendo:
\begin{equation*}
\begin{split}
H(X,Y)=
-\sum P(x,y)\log P(x)
-\sum P(x,y)\log P(y|x)
\end{split}
\end{equation*}
\sphinxAtStartPar
Primer término:
\begin{equation*}
\begin{split}
-\sum_x P(x)\log P(x)=H(X)
\end{split}
\end{equation*}
\sphinxAtStartPar
Segundo término:
\begin{equation*}
\begin{split}
-\sum P(x,y)\log P(y|x)=H(Y|X)
\end{split}
\end{equation*}
\sphinxAtStartPar
Concluimos:
\begin{equation*}
\begin{split}
H(X,Y)=H(X)+H(Y|X)
\end{split}
\end{equation*}

\bigskip\hrule\bigskip



\section{A.3 Simetría de la Información Mutua}
\label{\detokenize{05_appendices/proofs:a-3-simetria-de-la-informacion-mutua}}

\subsection{Teorema}
\label{\detokenize{05_appendices/proofs:id3}}\begin{equation*}
\begin{split}
I(X;Y)=I(Y;X)
\end{split}
\end{equation*}

\bigskip\hrule\bigskip



\subsection{Demostración}
\label{\detokenize{05_appendices/proofs:id4}}
\sphinxAtStartPar
Usando la identidad:
\begin{equation*}
\begin{split}
I(X;Y)=H(X)+H(Y)-H(X,Y)
\end{split}
\end{equation*}
\sphinxAtStartPar
y observando que:
\begin{equation*}
\begin{split}
H(X,Y)=H(Y,X)
\end{split}
\end{equation*}
\sphinxAtStartPar
la expresión es simétrica en \(X\) e \(Y\).

\sphinxAtStartPar
Por lo tanto:
\begin{equation*}
\begin{split}
I(X;Y)=I(Y;X)
\end{split}
\end{equation*}

\bigskip\hrule\bigskip



\section{A.4 Desigualdad de Gibbs}
\label{\detokenize{05_appendices/proofs:a-4-desigualdad-de-gibbs}}

\subsection{Teorema}
\label{\detokenize{05_appendices/proofs:id5}}
\sphinxAtStartPar
Para distribuciones \(P\) y \(Q\):
\begin{equation*}
\begin{split}
D_{KL}(P||Q)\ge0
\end{split}
\end{equation*}

\bigskip\hrule\bigskip



\subsection{Demostración}
\label{\detokenize{05_appendices/proofs:id6}}
\sphinxAtStartPar
Consideremos:
\begin{equation*}
\begin{split}
\log x \le x-1
\end{split}
\end{equation*}
\sphinxAtStartPar
para todo \(x>0\).

\sphinxAtStartPar
Sea:
\begin{equation*}
\begin{split}
x=\frac{Q(x)}{P(x)}
\end{split}
\end{equation*}
\sphinxAtStartPar
entonces:
\begin{equation*}
\begin{split}
-\log\frac{Q(x)}{P(x)}\ge1-\frac{Q(x)}{P(x)}
\end{split}
\end{equation*}
\sphinxAtStartPar
Multiplicando por \(P(x)\) y sumando:
\begin{equation*}
\begin{split}
\sum P(x)\log\frac{P(x)}{Q(x)}\ge0
\end{split}
\end{equation*}
\sphinxAtStartPar
lo que demuestra la no negatividad de la divergencia KL.


\bigskip\hrule\bigskip



\section{A.5 Regla de la Cadena para Información Mutua}
\label{\detokenize{05_appendices/proofs:a-5-regla-de-la-cadena-para-informacion-mutua}}

\subsection{Teorema}
\label{\detokenize{05_appendices/proofs:id7}}\begin{equation*}
\begin{split}
I(X;Y,Z)=I(X;Y)+I(X;Z|Y)
\end{split}
\end{equation*}

\bigskip\hrule\bigskip



\subsection{Demostración}
\label{\detokenize{05_appendices/proofs:id8}}
\sphinxAtStartPar
Partimos de:
\begin{equation*}
\begin{split}
I(X;Y,Z)=H(X)-H(X|Y,Z)
\end{split}
\end{equation*}
\sphinxAtStartPar
Usamos:
\begin{equation*}
\begin{split}
H(X|Y,Z)=H(X|Y)-I(X;Z|Y)
\end{split}
\end{equation*}
\sphinxAtStartPar
entonces:
\begin{equation*}
\begin{split}
I(X;Y,Z)=
H(X)-H(X|Y)+I(X;Z|Y)
\end{split}
\end{equation*}
\sphinxAtStartPar
pero:
\begin{equation*}
\begin{split}
H(X)-H(X|Y)=I(X;Y)
\end{split}
\end{equation*}
\sphinxAtStartPar
por lo que:
\begin{equation*}
\begin{split}
I(X;Y,Z)=I(X;Y)+I(X;Z|Y)
\end{split}
\end{equation*}

\bigskip\hrule\bigskip



\section{A.6 Desigualdad de Procesamiento de Datos}
\label{\detokenize{05_appendices/proofs:a-6-desigualdad-de-procesamiento-de-datos}}

\subsection{Teorema}
\label{\detokenize{05_appendices/proofs:id9}}
\sphinxAtStartPar
Si:
\begin{equation*}
\begin{split}
X \rightarrow Y \rightarrow Z
\end{split}
\end{equation*}
\sphinxAtStartPar
forma una cadena de Markov, entonces:
\begin{equation*}
\begin{split}
I(X;Z)\le I(X;Y)
\end{split}
\end{equation*}

\bigskip\hrule\bigskip



\subsection{Demostración (Esquema Formal)}
\label{\detokenize{05_appendices/proofs:demostracion-esquema-formal}}
\sphinxAtStartPar
Propiedad de Markov:
\begin{equation*}
\begin{split}
P(z|x,y)=P(z|y)
\end{split}
\end{equation*}
\sphinxAtStartPar
Se demuestra que:
\begin{equation*}
\begin{split}
I(X;Z)=I(X;Y)-I(X;Y|Z)+I(X;Z|Y)
\end{split}
\end{equation*}
\sphinxAtStartPar
pero:
\begin{equation*}
\begin{split}
I(X;Z|Y)=0
\end{split}
\end{equation*}
\sphinxAtStartPar
bajo la condición de Markov.

\sphinxAtStartPar
Dado que:
\begin{equation*}
\begin{split}
I(X;Y|Z)\ge0
\end{split}
\end{equation*}
\sphinxAtStartPar
se obtiene:
\begin{equation*}
\begin{split}
I(X;Z)\le I(X;Y)
\end{split}
\end{equation*}

\bigskip\hrule\bigskip



\section{A.7 Descomposición de Información Mutua Multivariada}
\label{\detokenize{05_appendices/proofs:a-7-descomposicion-de-informacion-mutua-multivariada}}

\subsection{Teorema}
\label{\detokenize{05_appendices/proofs:id10}}\begin{equation*}
\begin{split}
I(X;Y,Z,W)
=I(X;Y)+I(X;Z|Y)+I(X;W|Y,Z)
\end{split}
\end{equation*}

\bigskip\hrule\bigskip



\subsection{Demostración}
\label{\detokenize{05_appendices/proofs:id11}}
\sphinxAtStartPar
Aplicamos repetidamente la regla de la cadena:
\begin{equation*}
\begin{split}
I(X;Y,Z,W)=I(X;Y)+I(X;Z,W|Y)
\end{split}
\end{equation*}
\sphinxAtStartPar
Luego:
\begin{equation*}
\begin{split}
I(X;Z,W|Y)=I(X;Z|Y)+I(X;W|Y,Z)
\end{split}
\end{equation*}
\sphinxAtStartPar
Combinando:
\begin{equation*}
\begin{split}
I(X;Y,Z,W)=I(X;Y)+I(X;Z|Y)+I(X;W|Y,Z)
\end{split}
\end{equation*}

\bigskip\hrule\bigskip



\section{Conclusión del Apéndice}
\label{\detokenize{05_appendices/proofs:conclusion-del-apendice}}
\sphinxAtStartPar
Las demostraciones anteriores establecen los fundamentos formales para:
\begin{itemize}
\item {} 
\sphinxAtStartPar
teoría de la información clásica

\item {} 
\sphinxAtStartPar
modelos multivariados

\item {} 
\sphinxAtStartPar
aprendizaje profundo informacional

\item {} 
\sphinxAtStartPar
análisis estadístico de dependencias

\end{itemize}

\sphinxAtStartPar
Sirven como base matemática rigurosa para los capítulos teóricos y aplicaciones presentados en el libro.

\sphinxstepscope


\chapter{Tablas Avanzadas de Teoría de la Información}
\label{\detokenize{05_appendices/advanced_tables:tablas-avanzadas-de-teoria-de-la-informacion}}\label{\detokenize{05_appendices/advanced_tables::doc}}

\section{Introducción}
\label{\detokenize{05_appendices/advanced_tables:introduccion}}
\sphinxAtStartPar
Este apéndice reúne tablas de referencia rápida con definiciones, propiedades,
desigualdades fundamentales y relaciones estructurales entre medidas de la
teoría de la información.\\
El objetivo es proporcionar un compendio operativo útil para investigación,
docencia y aplicaciones en ciencia de datos e ingeniería.


\bigskip\hrule\bigskip



\section{B.1 Definiciones Fundamentales}
\label{\detokenize{05_appendices/advanced_tables:b-1-definiciones-fundamentales}}

\begin{savenotes}\sphinxattablestart
\centering
\begin{tabulary}{\linewidth}[t]{|T|T|T|}
\hline
\sphinxstyletheadfamily 
\sphinxAtStartPar
Medida
&\sphinxstyletheadfamily 
\sphinxAtStartPar
Definición Matemática
&\sphinxstyletheadfamily 
\sphinxAtStartPar
Interpretación
\\
\hline
\sphinxAtStartPar
Entropía
&
\sphinxAtStartPar
\(H(X)=-\sum_x P(x)\log P(x)\)
&
\sphinxAtStartPar
Incertidumbre promedio
\\
\hline
\sphinxAtStartPar
Entropía conjunta
&
\sphinxAtStartPar
\(H(X,Y)\)
&
\sphinxAtStartPar
Incertidumbre del sistema completo
\\
\hline
\sphinxAtStartPar
Entropía condicional
&
\sphinxAtStartPar
\$H(X
&
\sphinxAtStartPar
Y)\$
\\
\hline
\sphinxAtStartPar
Información mutua
&
\sphinxAtStartPar
\(I(X;Y)\)
&
\sphinxAtStartPar
Dependencia estadística
\\
\hline
\sphinxAtStartPar
MI condicional
&
\sphinxAtStartPar
\$I(X;Y
&
\sphinxAtStartPar
Z)\$
\\
\hline
\sphinxAtStartPar
Divergencia KL
&
\sphinxAtStartPar
\$D\_\{KL\}(P
&
\sphinxAtStartPar

\\
\hline
\end{tabulary}
\par
\sphinxattableend\end{savenotes}


\bigskip\hrule\bigskip



\section{B.2 Relaciones Fundamentales}
\label{\detokenize{05_appendices/advanced_tables:b-2-relaciones-fundamentales}}

\begin{savenotes}\sphinxattablestart
\centering
\begin{tabulary}{\linewidth}[t]{|T|T|}
\hline
\sphinxstyletheadfamily 
\sphinxAtStartPar
Relación
&\sphinxstyletheadfamily 
\sphinxAtStartPar
Expresión
\\
\hline
\sphinxAtStartPar
Regla de cadena
&
\sphinxAtStartPar
\$H(X,Y)=H(X)+H(Y
\\
\hline
\sphinxAtStartPar
MI por entropías
&
\sphinxAtStartPar
\(I(X;Y)=H(X)+H(Y)-H(X,Y)\)
\\
\hline
\sphinxAtStartPar
MI alternativa
&
\sphinxAtStartPar
\$I(X;Y)=H(X)\sphinxhyphen{}H(X
\\
\hline
\sphinxAtStartPar
MI condicional
&
\sphinxAtStartPar
\$I(X;Y
\\
\hline
\sphinxAtStartPar
Cadena MI
&
\sphinxAtStartPar
\$I(X;Y,Z)=I(X;Y)+I(X;Z
\\
\hline
\end{tabulary}
\par
\sphinxattableend\end{savenotes}


\bigskip\hrule\bigskip



\section{B.3 Desigualdades Clásicas}
\label{\detokenize{05_appendices/advanced_tables:b-3-desigualdades-clasicas}}

\begin{savenotes}\sphinxattablestart
\centering
\begin{tabulary}{\linewidth}[t]{|T|T|}
\hline
\sphinxstyletheadfamily 
\sphinxAtStartPar
Desigualdad
&\sphinxstyletheadfamily 
\sphinxAtStartPar
Expresión
\\
\hline
\sphinxAtStartPar
No negatividad KL
&
\sphinxAtStartPar
\$D\_\{KL\}(P
\\
\hline
\sphinxAtStartPar
No negatividad MI
&
\sphinxAtStartPar
\(I(X;Y)\ge0\)
\\
\hline
\sphinxAtStartPar
Subaditividad
&
\sphinxAtStartPar
\(H(X,Y)\le H(X)+H(Y)\)
\\
\hline
\sphinxAtStartPar
Monotonía
&
\sphinxAtStartPar
\$H(X
\\
\hline
\sphinxAtStartPar
Procesamiento de datos
&
\sphinxAtStartPar
\(I(X;Z)\le I(X;Y)\)
\\
\hline
\end{tabulary}
\par
\sphinxattableend\end{savenotes}


\bigskip\hrule\bigskip



\section{B.4 Propiedades Algebraicas}
\label{\detokenize{05_appendices/advanced_tables:b-4-propiedades-algebraicas}}

\begin{savenotes}\sphinxattablestart
\centering
\begin{tabulary}{\linewidth}[t]{|T|T|T|}
\hline
\sphinxstyletheadfamily 
\sphinxAtStartPar
Propiedad
&\sphinxstyletheadfamily 
\sphinxAtStartPar
Entropía
&\sphinxstyletheadfamily 
\sphinxAtStartPar
Información Mutua
\\
\hline
\sphinxAtStartPar
Simetría
&
\sphinxAtStartPar
\(H(X,Y)=H(Y,X)\)
&
\sphinxAtStartPar
\(I(X;Y)=I(Y;X)\)
\\
\hline
\sphinxAtStartPar
No negatividad
&
\sphinxAtStartPar
\(H(X)\ge0\)
&
\sphinxAtStartPar
\(I(X;Y)\ge0\)
\\
\hline
\sphinxAtStartPar
Invariancia
&
\sphinxAtStartPar
Bajo permutaciones
&
\sphinxAtStartPar
Bajo permutaciones
\\
\hline
\sphinxAtStartPar
Extensibilidad
&
\sphinxAtStartPar
Sí
&
\sphinxAtStartPar
Sí
\\
\hline
\end{tabulary}
\par
\sphinxattableend\end{savenotes}


\bigskip\hrule\bigskip



\section{B.5 Descomposición Multivariada}
\label{\detokenize{05_appendices/advanced_tables:b-5-descomposicion-multivariada}}

\begin{savenotes}\sphinxattablestart
\centering
\begin{tabulary}{\linewidth}[t]{|T|T|}
\hline
\sphinxstyletheadfamily 
\sphinxAtStartPar
Expansión
&\sphinxstyletheadfamily 
\sphinxAtStartPar
Expresión
\\
\hline
\sphinxAtStartPar
Entropía triple
&
\sphinxAtStartPar
\$H(X,Y,Z)=H(X)+H(Y
\\
\hline
\sphinxAtStartPar
MI triple
&
\sphinxAtStartPar
\$I(X;Y,Z)=I(X;Y)+I(X;Z
\\
\hline
\sphinxAtStartPar
MI cuádruple
&
\sphinxAtStartPar
\$I(X;Y,Z,W)=I(X;Y)+I(X;Z
\\
\hline
\end{tabulary}
\par
\sphinxattableend\end{savenotes}


\bigskip\hrule\bigskip



\section{B.6 Interpretaciones en Ciencia de Datos}
\label{\detokenize{05_appendices/advanced_tables:b-6-interpretaciones-en-ciencia-de-datos}}

\begin{savenotes}\sphinxattablestart
\centering
\begin{tabulary}{\linewidth}[t]{|T|T|}
\hline
\sphinxstyletheadfamily 
\sphinxAtStartPar
Medida
&\sphinxstyletheadfamily 
\sphinxAtStartPar
Uso típico
\\
\hline
\sphinxAtStartPar
Entropía
&
\sphinxAtStartPar
Medición de complejidad
\\
\hline
\sphinxAtStartPar
MI
&
\sphinxAtStartPar
Selección de características
\\
\hline
\sphinxAtStartPar
MI condicional
&
\sphinxAtStartPar
Eliminación de redundancia
\\
\hline
\sphinxAtStartPar
KL
&
\sphinxAtStartPar
Evaluación de modelos
\\
\hline
\sphinxAtStartPar
Entropía conjunta
&
\sphinxAtStartPar
Modelado multivariado
\\
\hline
\end{tabulary}
\par
\sphinxattableend\end{savenotes}


\bigskip\hrule\bigskip



\section{B.7 Propiedades en Aprendizaje Automático}
\label{\detokenize{05_appendices/advanced_tables:b-7-propiedades-en-aprendizaje-automatico}}

\begin{savenotes}\sphinxattablestart
\centering
\begin{tabulary}{\linewidth}[t]{|T|T|}
\hline
\sphinxstyletheadfamily 
\sphinxAtStartPar
Contexto
&\sphinxstyletheadfamily 
\sphinxAtStartPar
Uso de MI
\\
\hline
\sphinxAtStartPar
Feature Selection
&
\sphinxAtStartPar
Ranking informacional
\\
\hline
\sphinxAtStartPar
Clustering
&
\sphinxAtStartPar
Medida de similitud
\\
\hline
\sphinxAtStartPar
Deep Learning
&
\sphinxAtStartPar
Bottleneck informacional
\\
\hline
\sphinxAtStartPar
Representaciones latentes
&
\sphinxAtStartPar
Maximización MI
\\
\hline
\sphinxAtStartPar
Regularización
&
\sphinxAtStartPar
Penalización KL
\\
\hline
\end{tabulary}
\par
\sphinxattableend\end{savenotes}


\bigskip\hrule\bigskip



\section{B.8 Propiedades de Independencia}
\label{\detokenize{05_appendices/advanced_tables:b-8-propiedades-de-independencia}}

\begin{savenotes}\sphinxattablestart
\centering
\begin{tabulary}{\linewidth}[t]{|T|T|}
\hline
\sphinxstyletheadfamily 
\sphinxAtStartPar
Condición
&\sphinxstyletheadfamily 
\sphinxAtStartPar
Consecuencia
\\
\hline
\sphinxAtStartPar
Independencia
&
\sphinxAtStartPar
\(I(X;Y)=0\)
\\
\hline
\sphinxAtStartPar
Independencia condicional
&
\sphinxAtStartPar
\$I(X;Y
\\
\hline
\sphinxAtStartPar
Cadena de Markov
&
\sphinxAtStartPar
Procesamiento de datos
\\
\hline
\sphinxAtStartPar
Factorización
&
\sphinxAtStartPar
\(P(x,y)=P(x)P(y)\)
\\
\hline
\end{tabulary}
\par
\sphinxattableend\end{savenotes}


\bigskip\hrule\bigskip



\section{B.9 Propiedades Diferenciales (Continuas)}
\label{\detokenize{05_appendices/advanced_tables:b-9-propiedades-diferenciales-continuas}}

\begin{savenotes}\sphinxattablestart
\centering
\begin{tabulary}{\linewidth}[t]{|T|T|}
\hline
\sphinxstyletheadfamily 
\sphinxAtStartPar
Medida
&\sphinxstyletheadfamily 
\sphinxAtStartPar
Forma Continua
\\
\hline
\sphinxAtStartPar
Entropía diferencial
&
\sphinxAtStartPar
\(h(X)=-\int p(x)\log p(x)\,dx\)
\\
\hline
\sphinxAtStartPar
MI continua
&
\sphinxAtStartPar
Integral doble
\\
\hline
\sphinxAtStartPar
KL continua
&
\sphinxAtStartPar
Integral sobre densidades
\\
\hline
\end{tabulary}
\par
\sphinxattableend\end{savenotes}


\bigskip\hrule\bigskip



\section{B.10 Complejidad Computacional (Estimación MI)}
\label{\detokenize{05_appendices/advanced_tables:b-10-complejidad-computacional-estimacion-mi}}

\begin{savenotes}\sphinxattablestart
\centering
\begin{tabulary}{\linewidth}[t]{|T|T|}
\hline
\sphinxstyletheadfamily 
\sphinxAtStartPar
Método
&\sphinxstyletheadfamily 
\sphinxAtStartPar
Complejidad Aproximada
\\
\hline
\sphinxAtStartPar
Histogramas
&
\sphinxAtStartPar
\(O(n)\)
\\
\hline
\sphinxAtStartPar
KDE
&
\sphinxAtStartPar
\(O(n^2)\)
\\
\hline
\sphinxAtStartPar
kNN
&
\sphinxAtStartPar
\(O(n\log n)\)
\\
\hline
\sphinxAtStartPar
MINE
&
\sphinxAtStartPar
Dependiente del modelo
\\
\hline
\end{tabulary}
\par
\sphinxattableend\end{savenotes}


\bigskip\hrule\bigskip



\section{B.11 Resumen de Relaciones Clave}
\label{\detokenize{05_appendices/advanced_tables:b-11-resumen-de-relaciones-clave}}\begin{equation*}
\begin{split}
\begin{aligned}
I(X;Y) &= H(X)-H(X|Y) \\
I(X;Y) &= D_{KL}(P(X,Y)||P(X)P(Y)) \\
H(X,Y) &= H(X)+H(Y|X) \\
I(X;Y|Z) &= H(X|Z)-H(X|Y,Z)
\end{aligned}
\end{split}
\end{equation*}

\bigskip\hrule\bigskip



\section{Conclusión del Apéndice}
\label{\detokenize{05_appendices/advanced_tables:conclusion-del-apendice}}
\sphinxAtStartPar
Las tablas anteriores sintetizan las relaciones estructurales más importantes
de la teoría de la información y constituyen una referencia rápida para:
\begin{itemize}
\item {} 
\sphinxAtStartPar
desarrollo teórico

\item {} 
\sphinxAtStartPar
análisis estadístico

\item {} 
\sphinxAtStartPar
modelamiento multivariado

\item {} 
\sphinxAtStartPar
aplicaciones en ciencia de datos y aprendizaje automático

\end{itemize}

\sphinxstepscope


\chapter{Agradecimientos}
\label{\detokenize{05_appendices/acknowledgments:agradecimientos}}\label{\detokenize{05_appendices/acknowledgments::doc}}
\sphinxAtStartPar
:class: unnumbered

\sphinxAtStartPar
El autor expresa su sincero agradecimiento a la comunidad académica
de teoría de la información, aprendizaje automático y ciencia de datos,
cuyas contribuciones han permitido el desarrollo conceptual presentado
en este libro.

\sphinxAtStartPar
Se reconoce especialmente la influencia de los trabajos clásicos de
Claude Shannon, Thomas Cover, David MacKay y Kevin Murphy,
cuyas obras constituyen pilares fundamentales en la disciplina.

\sphinxAtStartPar
Asimismo, se agradece a la comunidad científica abierta por el
desarrollo de herramientas computacionales como Python, Jupyter
y el ecosistema científico que posibilita la reproducibilidad
del conocimiento.

\sphinxstepscope


\chapter{Bibliografía}
\label{\detokenize{05_appendices/references:bibliografia}}\label{\detokenize{05_appendices/references::doc}}
\sphinxAtStartPar
:class: unnumbered

\begin{sphinxthebibliography}{Bis06}
\bibitem[Bis06]{05_appendices/references:id4}
\sphinxAtStartPar
Christopher Bishop. \sphinxstyleemphasis{Pattern Recognition and Machine Learning}. Springer, 2006.
\bibitem[CT06]{05_appendices/references:id2}
\sphinxAtStartPar
Thomas Cover and Joy Thomas. \sphinxstyleemphasis{Elements of Information Theory}. Wiley, 2006.
\bibitem[Csi11]{05_appendices/references:id10}
\sphinxAtStartPar
Imre Csiszár. \sphinxstyleemphasis{Information Theory and Statistics}. Springer, 2011.
\bibitem[GBC16]{05_appendices/references:id8}
\sphinxAtStartPar
Ian Goodfellow, Yoshua Bengio, and Aaron Courville. \sphinxstyleemphasis{Deep Learning}. MIT Press, 2016.
\bibitem[Gra11]{05_appendices/references:id12}
\sphinxAtStartPar
Robert Gray. \sphinxstyleemphasis{Entropy and Information Theory}. Springer, 2011.
\bibitem[KSG04]{05_appendices/references:id7}
\sphinxAtStartPar
Alexander Kraskov, Harald Stögbauer, and Peter Grassberger. Estimating mutual information. \sphinxstyleemphasis{Physical Review E}, 69(6):066138, 2004.
\bibitem[KL51]{05_appendices/references:id15}
\sphinxAtStartPar
Solomon Kullback and Richard Leibler. On information and sufficiency. \sphinxstyleemphasis{Annals of Mathematical Statistics}, 1951.
\bibitem[Mac03]{05_appendices/references:id3}
\sphinxAtStartPar
David MacKay. \sphinxstyleemphasis{Information Theory, Inference, and Learning Algorithms}. Cambridge University Press, 2003.
\bibitem[Mur12]{05_appendices/references:id5}
\sphinxAtStartPar
Kevin Murphy. \sphinxstyleemphasis{Machine Learning: A Probabilistic Perspective}. MIT Press, 2012.
\bibitem[Pea09]{05_appendices/references:id9}
\sphinxAtStartPar
Judea Pearl. \sphinxstyleemphasis{Causality}. Cambridge, 2009.
\bibitem[Pen05]{05_appendices/references:id13}
\sphinxAtStartPar
Hanchuan Peng. Feature selection based on mutual information. \sphinxstyleemphasis{IEEE TPAMI}, 2005.
\bibitem[Sha48]{05_appendices/references:id6}
\sphinxAtStartPar
Claude E. Shannon. A mathematical theory of communication. \sphinxstyleemphasis{Bell System Technical Journal}, 27(3):379–423, 1948.
\bibitem[Tis00]{05_appendices/references:id14}
\sphinxAtStartPar
Naftali Tishby. The information bottleneck method. \sphinxstyleemphasis{arXiv preprint physics/0004057}, 2000.
\bibitem[Yeu08]{05_appendices/references:id11}
\sphinxAtStartPar
Raymond Yeung. \sphinxstyleemphasis{Information Theory and Network Coding}. Springer, 2008.
\end{sphinxthebibliography}







\renewcommand{\indexname}{Index}
\printindex
\end{document}